\documentclass[final,3p,times,fleqn]{elsarticle}

\usepackage{textgreek}
\usepackage{stackengine}

\newcommand\textdot[1]{\stackon[1pt]{\csname text#1\endcsname}{.}}
\newcommand\textdots[1]{\stackon[1pt]{\csname text#1\endcsname}{..}}

\parindent 1 pc

\usepackage{graphicx}
\usepackage{epsfig}

\usepackage{amsmath}
\usepackage{amssymb}
\usepackage{amsthm}
\usepackage{mathrsfs}
\usepackage{amsfonts}
\usepackage{mathtools}

\usepackage{setspace}
\usepackage{natbib}
\usepackage{multirow}
\usepackage{booktabs}
\usepackage{xcolor}
\definecolor{revcolor}{RGB}{0,0,200}
\newcommand{\rev}[1]{\textcolor{revcolor}{#1}}
\usepackage{nomencl}
\makenomenclature
\usepackage[running]{lineno}
\usepackage{makecell}
\usepackage{subcaption}
\usepackage{wrapfig}

\biboptions{sort&compress}

\usepackage[colorlinks=true,
            linkcolor=blue,
            urlcolor=blue,
            citecolor=blue,
            anchorcolor=blue]{hyperref}

\linespread{1.3}
\makeatletter

\parindent 10 pt
\setlength{\topmargin}{-22.0 mm}
\setlength{\headheight}{4.0 mm}
\setlength{\headsep}{10.0 mm}
\setlength{\topskip}{0.00 mm}
\setlength{\textheight}{252.0 mm}
\setlength{\textwidth}{170.0 mm}
\setlength{\oddsidemargin}{-6.0 mm}
\setlength{\evensidemargin}{-3.0 mm}

\newcommand{\Lop}{\mathcal{L}}

\input epsf

\newcounter{saveeqn}
\newcommand{\alpheqn}{\setcounter{saveeqn}{\value{equation}}%
\stepcounter{saveeqn}\setcounter{equation}{1}%
\renewcommand{\theequation}{\mbox{\arabic{saveeqn}-\alph{equation}}}}

\newcommand{\reseteqn}{\setcounter{equation}{\value{saveeqn}}%
\renewcommand{\theequation}{\arabic{equation}}}

\newcommand{\R}{{\mbox{Re}}}
\newcommand{\W}{{\mbox{We}}}
\newcommand{\Fr}{{\mbox{Fr}}}
\newcommand{\Fi}{{\mbox{Fi}}}
\newcommand{\M}{{\mbox{M}}}
\newcommand{\E}{{\mbox{E}}}

\begin{document}

\begin{frontmatter}

\title{Finite Element Analysis of Second-Order Strain-Gradient Elasticity with Ellipsoidal Nonlocal Characteristic Lengths}


\author{
Vipin Gupta$^{1}$,
Sunita Kumawat$^{2}$,
Arun Kumar Yadav$^{3}$,
Taizo Maruyama$^{4}$,
Bandar Almohsen$^{5,*}$
\\
${}^{1}$Department of Mathematics, Gurugram University, Gurugram, India\\
${}^{2}$Department of Mathematics, Amrita Vishwa Vidyapeetham, Bengaluru Campus, Bengaluru-560035, India\\
${}^{3}$School of Basic and Applied Sciences, K.R. Mangalam University, Gurugram, India\\
${}^{4}$Department of Civil and Environmental Engineering, Institute of Science Tokyo, 2-12-1-W8-22, Ookayama, Meguro-ku, Tokyo 152-8550, Japan\\
${}^{5}$Department of Mathematics, College of Science, King Saud University, P.O. Box 2455, Riyadh 11451, Saudi Arabia\\
Email: $^{1}$vipin@gurugramuniversity.ac.in\\
$^{2}$sunita14kumawat@gmail.com\\
$^{3}$arunyadav0307@gmail.com\\
$^{4}$maruyama.t.45ef@m.isct.ac.jp\\
$^{5}$balmohsen@ksu.edu.sa\\
${}^{*}$Corresponding author
}
%====================  ABSTRACT STARTS BELOW  ==================



\begin{abstract}
\rev{This study addresses how direction-dependent microstructural interactions can be derived from a geometrically defined nonlocal domain and incorporated into a unified finite element framework for anisotropic strain-gradient elasticity. A three-dimensional ellipsoidal nonlocal averaging procedure is used to derive an anisotropic characteristic-length tensor $L_{mn}=(AA^T)_{mn}$ from the second moments of the averaging domain, giving a geometric interpretation analogous to fabric tensor. A positive-definite second-order strain-gradient energy $W=\frac12 C_{ijkl}\varepsilon_{ij}\varepsilon_{kl}+\frac1{20}L_{mn}C_{ijkl}\varepsilon_{ij,m}\varepsilon_{kl,n}$ is formulated and implemented with $C^1$-continuous cubic Hermite (1D), Bogner-Fox-Schmit bicubic Hermite (2D) and tricubic Hermite (3D, 24 DOF/node) elements using a corrected consistent mass formulation. For $l_x=0.20$ m and $l_y=l_z=0.02$ m the orientation sweep shows monotonic variation with $u(\theta)/u(0^\circ)=1.00\to1.039$ and $f_1=186.26\to173.06$ Hz. With $f_{iso}=178.72$ Hz for $l_{ref}=0.05$ m (177.0 Hz for volume-equivalent $l_{iso}=0.0431$ m) the normalized frequency becomes $f_1/f_{iso}=1.054\to0.979$. The scaled condition number varies as $\kappa=2.7\times10^5\to10.3\times10^5$ at $u_{ref}=4.7876$ $\mu$m. All cases remain above the classical bound $f_1>171.61$ Hz. Axial length dominates the transverse contribution, giving $5.8\%$ reduction versus $0.5\%$. With $DSI=(u_{ref}-u)/u_{ref}$ defined against the volume-equivalent $l_{iso}=(l_xl_yl_z)^{1/3}$ recomputed per configuration, the index is predominantly positive from $0\%$ to $+29.6\%$, and only small negative values of $-0.1\%$ to $-0.9\%$ appear at high orientation angles for $AR\ge2$. This variation enables tailoring of stiffness and frequency through controlled anisotropy and orientation.}
\end{abstract}

\begin{keyword}
Strain-gradient elasticity; anisotropic nonlocality; ellipsoidal averaging neighbourhood; characteristic length tensor; \(C^1\)-continuous finite elements; size-dependent stiffening.
\end{keyword}

\end{frontmatter}

\section{Introduction}
Classical elasticity provides an efficient continuum description when the characteristic structural dimension is much larger than the underlying microstructural scale. Its constitutive relation is local: stress depends on strain at the same material point, and no intrinsic length parameter appears in the governing equations \cite{Love1927}.



However, it is restrictive for composites with microstructures, thin films, additively manufactured lattices and mechanical metamaterials, where the presence of pores, fibres, grains, inclusions and/or architected unit cells affect the apparent stiffness, vibration response, wave dispersion and boundary-layer behaviour \cite{Jiao2023,Sarar2024}. In these materials, a continuum model thus needs to contain characteristic lengths; these lengths should be able to take into account the directionally organized microstructure.


This is overcome by generalized continuum theories, which extend the kinematics or constitutive response. Theories of couple-stress and higher order elasticity were developed to incorporate measures of curvature, moment stresses, strain gradients, and higher order boundary conditions \cite{Toupin1962,MindlinTiersten1962,Mindlin1964,Mindlin1965}. Eringen-type nonlocal elasticity, on the other hand, viewed stress as an average response over a neighbourhood of material points \cite{EringenEdelen1972,Eringen1972Polar,Eringen1983} in parallel. The two are closely related: nonlocal averaging gives a physical interpretation to long-range interactions, while gradient approximations transform weak nonlocality into differential equations that can be solved numerically on a boundary value problem. In practice, however, they are not very useful because the general higher order constitutive tensors are too complicated, the nonlocal kernels are ambiguous, and consistent boundary conditions are hard to enforce on finite domains.




% This scale-free assumption becomes restrictive for microstructured composites, thin films, additively manufactured lattices, and mechanical metamaterials, where pores, fibres, grains, inclusions, or architected unit cells influence the apparent stiffness, vibration response, wave dispersion, and boundary-layer behaviour \cite{Jiao2023,Sarar2024}. For such materials, a continuum model must incorporate characteristic lengths, and these lengths should be capable of reflecting directional microstructural organization.

% Generalized continuum theories address this limitation by enriching the kinematics or constitutive response. Couple-stress and higher-order elasticity theories introduced curvature-type measures, moment stresses, strain gradients, and associated higher-order boundary conditions \cite{Toupin1962,MindlinTiersten1962,Mindlin1964,Mindlin1965}. In parallel, Eringen-type nonlocal elasticity described stress as a spatially averaged response over a neighbourhood of material points \cite{EringenEdelen1972,Eringen1972Polar,Eringen1983}. These two viewpoints are closely related: nonlocal averaging provides a physical interpretation of long-range interactions, whereas gradient approximations convert weak nonlocality into differential equations suitable for boundary-value problems. Their practical use, however, is limited by the complexity of general higher-order constitutive tensors, the ambiguity of nonlocal kernels, and the difficulty of imposing consistent boundary conditions in finite domains.



Within this framework, the simplified strain-gradient elasticity, notably the Aifantis form, provides a numerically more efficient alternative. These models regularize localization, eliminate singularities, and predict size-dependent responses by inclusion of a small internal length parameter \cite{Aifantis1992,RuAifantis1993,AltanAifantis1997}. 
Experimental and computational studies have confirmed the relevance of strain-gradient effects in microscale deformation and dynamic problems \cite{Lam2003,AskesAifantis2011}. Discrete-to-continuum and homogenization-based arguments further show that gradient terms may emerge naturally from lattice or representative-volume descriptions \cite{MetrikineAskes2002,Gitman2005}. Yet most simplified models use a scalar length scale, implicitly assuming that nonlocal interactions are isotropic. This assumption is inadequate for layered media, fibre-reinforced solids, anisotropic crystals, and architected lattices, where the interaction range may differ substantially from one direction to another.

Anisotropic gradient elasticity has therefore become an important extension of scalar-length models. Tensorial characteristic lengths allow the weak nonlocal response to depend on direction, and this idea has been developed through separable anisotropic gradient, internal-length tensor formulations, and symmetry classifications of three-dimensional strain-gradient elasticity \cite{LazarPo2015,Polizzotto2018,Auffray2019}. These studies establish that the internal length is not necessarily a scalar material constant; it may be a second-order tensor carrying information about microstructural orientation. However, in many formulations this tensor is introduced phenomenologically or through homogenized coefficients. The link between a geometrically defined nonlocal averaging domain and the resulting anisotropic length tensor remains less explicit, although such a link is essential for parameter interpretation and model calibration.

Recent work reinforces the same need from different perspectives. Physically based nonlocal strain-gradient models have related kernel functions and intrinsic lengths to polymer-network descriptors \cite{Jiang2023}, while second-order homogenization of three-dimensional lattice materials has shown that strain-gradient continua can reproduce size effects that classical homogenization misses \cite{Molavitabrizi2023}.
\rev{Recent developments in non-classical continuum theories have expanded applications to smart structures and multiphysics. Nonlocal strain gradient models have been applied to wave propagation in functionally graded materials \cite{PhungVan2024AMM128} and to size-dependent bending of porous FG beams \cite{PhungVan2024AMM127}. Other applications include buckling and vibration of composite laminates with nonlocal effects \cite{PhungVan2024CS335} and isogeometric analysis of gradient plates \cite{PhungVan2023EC39}. Boundary element formulations for nonlocal elasticity \cite{Schwartz2012EABE36} and strain-gradient effects in thin-walled structures \cite{PhungVan2024TWS198} have also been reported. These studies confirm that internal length parameters significantly modify stiffness, vibration and wave dispersion, motivating anisotropic extensions where length itself becomes tensorial.} Hierarchical lattice studies further indicate that higher-order stiffness tensors improve static and dynamic predictions in architected solids \cite{YangLiuXia2024}. At the defect scale, nonlocal simplified strain-gradient elasticity has clarified the joint role of Eringen-type and Aifantis-type lengths in regularizing dislocation fields \cite{Lazar2024}. Related nonlocal multiphysics studies, including the recent work of Gupta and co-workers on piezo-thermoelastic, piezo-semiconductor, voided, and fractional Moore--Gibson--Thompson media, show that nonlocal parameters can significantly modify wave speed, attenuation, vibration response, and particle motion in smart small-scale systems \cite{Gupta2024Impact,Gupta2024Vibrational,Gupta2024Rayleigh,Gupta2025Piezo}.



Although these developments broaden the scope of nonlocal and gradient modelling, they also highlight a remaining limitation: most formulations are either analytical and geometry-specific, homogenization-dependent, or not directly connected to a unified finite element framework for anisotropic finite-domain problems. Finally, the numerical treatment of strain-gradient elasticity represents an additional challenge. In a primal (displacement-based) Galerkin formulation, fourth-order governing equations necessitate $C^{1}$-continuous approximation. The mixed and staggered finite element methods relax the continuity requirement but add new fields, stabilization choices or a saddle-point structure \cite{AmanatidouAravas2002,AskesMorataAifantis2008,PhunpengBaiz2015}.

The direct $C^{1}$ formulations preserve the variational structure, and when appropriate interpolation functions exist, they are also convenient to apply for general gradient elasticity. For example, such approximations can be achieved in a classical sense via cubic Hermite interpolation, Bogner--Fox--Schmit elements and isoparametric Hermite constructions \cite{BognerFoxSchmit1966,PeteraPittman1994}. Recently, finite element developments for elasticity with microstructure, three-dimensional $C^{1}$ gradient elasticity and associated verification benchmarks \cite{Zervos2008,Papanicolopulos2009}, as well as virtual element formulations \cite{Yang2021,WriggersHudobivnik2023} have elucidated the computational requirements of higher-order continua. A unified one-, two-, and three-dimensional $C^{1}$ implementation coupled with anisotropic characteristic lengths derived from an explicit ellipsoidal nonlocal averaging domain is still not commonly available.

To address the above limitations, the present study derives anisotropic characteristic lengths directly from ellipsoidal nonlocal strain averaging, thereby providing a geometric interpretation of the anisotropic length tensor. Based on this framework, a positive-definite second-order strain-gradient formulation is developed and incorporated into a unified $C^{1}$-continuous finite element framework employing cubic Hermite, Bogner--Fox--Schmit, and tensor-product tricubic Hermite elements for one-, two-, and three-dimensional analyses.

The main contributions of this work are:
\begin{itemize}
    \item \rev{Ellipsoidal nonlocal averaging is used to derive direction-dependent characteristic lengths. The resulting tensor $L_{mn}=(AA^T)_{mn}$ represents second moments of the averaging domain and provides a geometric interpretation analogous to fabric tensors in homogenization literature \cite{LazarPo2015,Polizzotto2018,Auffray2019}.}
    \item \rev{A positive-definite anisotropic second-order strain-gradient energy $W=1/2 C_{ijkl}\varepsilon_{ij}\varepsilon_{kl}+1/20 L_{mn}C_{ijkl}\varepsilon_{ij,m}\varepsilon_{kl,n}$ is formulated together with constitutive relations, governing equations and higher-order boundary conditions. After variation the $1/20$ factor becomes $1/10$ in the stress, which gives a positive stiffness correction in the wavenumber domain.}
    \item \rev{A unified $C^{1}$-continuous finite element framework is developed for 1D, 2D and 3D. It uses cubic Hermite, Bogner--Fox--Schmit bicubic Hermite and tensor-product tricubic Hermite elements with 24 DOF per node in 3D, and diagonal scaling keeps the condition number $\kappa<10^6$.}
    \item \rev{Orientation-dependent tuning is governed by $L_{11}(\theta)=l_x^2\cos^2\theta+l_y^2\sin^2\theta$ and the off-diagonal term $L_{12}=(l_y^2-l_x^2)\cos\theta\sin\theta$ that couples $B_{,x}^T C B_{,y}$. With the corrected mass $M_{scaled}=S_{scale}M_{ff}S_{scale}$, the response varies monotonically with orientation: $f_1$ decreases from $186.26$ to $173.06$ Hz and $u/u_0$ increases from $1.00$ to $1.039$. The normalized frequency becomes $f_1/f_{iso}=1.054\to0.979$ with no minimum at $45^\circ$. All cases remain above the classical bound $f_1>171.61$ Hz.}
    \item \rev{A bivariate design-space map in $AR=l_x/l_y\in[1,20]$ and $\theta\in[0^\circ,90^\circ]$ is constructed together with the Directional Stiffening Index $DSI=(u_{ref}-u)/u_{ref}$. The index quantifies displacement reduction relative to the volume-equivalent isotropic reference $l_{iso}=(l_xl_yl_z)^{1/3}$. The map reveals three regions: stiffening-optimal at high AR and low $\theta$, frequency-tuning at high AR near $\theta\approx45^\circ$, and insensitive when $AR\lesssim2$ or $\theta\approx90^\circ$.}
\end{itemize}



% -------------------------% ===================================================================
\section{Formulation of the Problem}

\subsection{Nonlocal Strain Averaging}

Let $\Omega\subset\mathbb{R}^{3}$ denote the material body and let
$\mathbf{x}\in\Omega$ be an interior point such that the nonlocal
neighbourhood introduced below remains contained in $\Omega$. The
nonlocal strain is introduced as a volume average of the classical strain
over a finite neighbourhood of $\mathbf{x}$, following the weakly
nonlocal interpretation of Eringen-type continuum theories
\cite{EringenEdelen1972,Eringen1972Polar,Eringen1983}:
\begin{equation}\label{eq:nonlocal_strain}
    \tilde{\varepsilon}_{mn}(\mathbf{x})
    =
    \frac{1}{|\mathcal{V}|}
    \int_{\mathcal{V}}
    p(\boldsymbol{\xi})\,
    \varepsilon_{mn}(\mathbf{x}+\boldsymbol{\xi})\,dV_{\xi}.
\end{equation}
Here $\varepsilon_{mn}$ is the classical infinitesimal strain tensor,
$\boldsymbol{\xi}$ is the relative position vector inside the averaging
neighbourhood, and $\mathcal{V}$ is the nonlocal interaction domain.
The kernel $p(\boldsymbol{\xi})$ is assumed to satisfy the normalization
condition
\begin{equation}\label{eq:kernel_normalization}
    \frac{1}{|\mathcal{V}|}
    \int_{\mathcal{V}}p(\boldsymbol{\xi})\,dV_{\xi}
    =
    1 .
\end{equation}
In the present formulation, a uniform kernel is adopted,
$p(\boldsymbol{\xi})=1$. This choice isolates the influence of the
anisotropic averaging geometry and leads to closed-form moment
relations.

The nonlocal neighbourhood is taken as an ellipsoid with semi-axes
$l_1,l_2,l_3>0$ aligned with the Cartesian coordinate axes:
\begin{equation}\label{eq:ellipsoid_domain}
\mathcal{V}
=
\left\{
\boldsymbol{\xi}\in\mathbb{R}^{3}
\;\middle|\;
\frac{\xi_1^2}{l_1^2}
+
\frac{\xi_2^2}{l_2^2}
+
\frac{\xi_3^2}{l_3^2}
\leq 1
\right\},
\qquad
|\mathcal{V}|=\frac{4\pi}{3}\,l_1l_2l_3 .
\end{equation}
The assumption that $\mathcal{V}$ is fully contained in $\Omega$
neglects boundary-induced truncation of the nonlocal neighbourhood.
The present derivation should therefore be interpreted as a bulk
weakly nonlocal approximation.
\rev{Near physical boundaries the averaging domain $\mathcal{V}(\mathbf{x})\cap\Omega$ is truncated. The first moment then becomes $\int\xi_i\neq0$ and the second moment becomes position-dependent, $L_{mn}(\mathbf{x})=1/|\mathcal{V}\cap\Omega|\int_{\mathcal{V}\cap\Omega}\xi_m\xi_n dV$. This correction introduces a boundary layer in the characteristic-length tensor itself and an additional first-gradient term. In the weakly nonlocal regime $l_{\max}\ll L_{\varepsilon}$ and for interior points with $l_{\max}\ll\text{dist}(\mathbf{x},\partial\Omega)$, the effect is second-order and is neglected. The natural double-traction condition $q_{ijm}\nu_m\nu_j=0$ therefore remains first-order consistent. A full surface-corrected formulation with $L_{mn}(\mathbf{x})$ and Mindlin-Toupin surface divergence terms lies beyond the present phenomenological Aifantis-type model and is noted as future work for thin structures where $l_{\max}\sim h$. For such thin structures the required modifications are as follows. First, $L_{mn}(\mathbf{x})$ is computed analytically for the given geometry. Second, the energy becomes spatially varying, $W(\mathbf{x})=1/2C\varepsilon\varepsilon+1/20 L_{mn}(\mathbf{x})C\varepsilon_{,m}\varepsilon_{,n}$. Third, extra terms proportional to $L_{mn,j}\varepsilon_{,mn}$ and $L_{mn,jj}\varepsilon$ appear in the governing equation. Finally, when $h\sim l_{\max}$ the full integral model $\tilde{\varepsilon}(\mathbf{x})=1/|\mathcal{V}\cap\Omega|\int_{\mathcal{V}\cap\Omega}\varepsilon(\mathbf{x}+\xi)dV$ is used without Taylor expansion.}

\begin{figure}[htbp]
\centering
\includegraphics[width=0.80\textwidth]{Fig1_Ellipsoidal_Neighbourhood.png}
\caption{Ellipsoidal nonlocal averaging neighbourhood and directional
characteristic lengths used in the anisotropic strain-gradient formulation.}
\label{fig:ellipsoidal_neighbourhood}
\end{figure}

The ellipsoid is mapped from the unit ball by
\[
    \boldsymbol{\xi}=A\mathbf{u},
    \qquad
    A=\operatorname{diag}(l_1,l_2,l_3),
    \qquad
    |\mathbf{u}|\leq 1 .
\]
Consequently,
\begin{equation}\label{eq:AAT}
    AA^T
    =
    \operatorname{diag}(l_1^2,l_2^2,l_3^2),
    \qquad
    (AA^T)_{ij}=l_i^2\delta_{ij}
    \quad \text{no sum on } i .
\end{equation}
The tensor $AA^T$ represents the directional second moment of the
averaging geometry and acts as an anisotropic characteristic-length
tensor. This interpretation is consistent with tensorial and ellipsoidal
descriptions of internal length in anisotropic gradient elasticity
\cite{LazarPo2015,Polizzotto2018,Auffray2019}. The present development
is restricted to an axis-aligned ellipsoidal neighbourhood, for which
$AA^T$ is diagonal. More generally, for a rotated ellipsoidal
neighbourhood, the characteristic-length tensor becomes
$QAA^TQ^T$, where $Q$ is an orthogonal rotation tensor.

Assuming that $\varepsilon_{mn}$ is sufficiently smooth in a
neighbourhood of $\mathbf{x}$, Taylor expansion gives
\begin{align}\label{eq:taylor}
\varepsilon_{mn}(\mathbf{x}+\boldsymbol{\xi})
&=
\varepsilon_{mn}(\mathbf{x})
+
\xi_i\,\varepsilon_{mn,i}
+
\frac{1}{2}\xi_i\xi_j\,\varepsilon_{mn,ij}
+
\frac{1}{6}\xi_i\xi_j\xi_k\,\varepsilon_{mn,ijk}
+
O(|\boldsymbol{\xi}|^4),
\end{align}
where a comma denotes spatial differentiation and repeated spatial
indices are summed. Since the ellipsoidal domain is centrosymmetric and
the adopted kernel is uniform, the first moment and all odd-order
moments vanish after integration. The required moment identities are
\begin{equation}\label{eq:moment_identities}
    \frac{1}{|\mathcal{V}|}
    \int_{\mathcal{V}}dV_{\xi}=1,
    \qquad
    \frac{1}{|\mathcal{V}|}
    \int_{\mathcal{V}}\xi_i\,dV_{\xi}=0,
    \qquad
    \frac{1}{|\mathcal{V}|}
    \int_{\mathcal{V}}\xi_i\xi_j\,dV_{\xi}
    =
    \frac{1}{5}(AA^T)_{ij}.
\end{equation}
The last identity follows from the transformation to the unit ball,
for which
\[
    \int_{|\mathbf{u}|\leq 1}u_i u_j\,dV_u
    =
    \frac{4\pi}{15}\delta_{ij}.
\]

Substitution of Eq.~\eqref{eq:taylor} into
Eq.~\eqref{eq:nonlocal_strain}, together with the moment identities in
Eq.~\eqref{eq:moment_identities}, gives the second-order weakly
nonlocal strain approximation
\begin{equation}\label{eq:nonlocal_strain_result}
\tilde{\varepsilon}_{mn}(\mathbf{x})
=
\varepsilon_{mn}(\mathbf{x})
+
\frac{1}{10}
(AA^T)_{ij}
\frac{\partial^2\varepsilon_{mn}}
     {\partial x_i\partial x_j}
+
O(l_{\max}^{4}\nabla^{4}\varepsilon_{mn}) ,
\end{equation}
where
\[
    l_{\max}=\max\{l_1,l_2,l_3\}.
\]
The coefficient $1/10$ is the product of the Taylor coefficient $1/2$
and the normalized second moment $1/5$ of the three-dimensional
ellipsoidal averaging domain. The approximation is valid in the weakly
nonlocal regime,
\[
    l_{\max}\ll L_{\varepsilon},
\]
where $L_{\varepsilon}$ denotes the characteristic length over which the
strain field varies appreciably. This Taylor-based reduction is the
standard connection between integral nonlocal averaging and
strain-gradient elasticity
\cite{Aifantis1992,AskesAifantis2011}.

Defining the anisotropic second-order differential operator
\begin{equation}\label{eq:Lop}
\mathcal{L}_{\ell}
\equiv
(AA^T)_{ij}
\frac{\partial^2}{\partial x_i\partial x_j}
=
l_1^2\frac{\partial^2}{\partial x_1^2}
+
l_2^2\frac{\partial^2}{\partial x_2^2}
+
l_3^2\frac{\partial^2}{\partial x_3^2},
\end{equation}
Eq.~\eqref{eq:nonlocal_strain_result} may be written compactly as
\begin{equation}\label{eq:nonlocal_strain_compact}
\tilde{\varepsilon}_{mn}(\mathbf{x})
=
\left(
1+\frac{1}{10}\mathcal{L}_{\ell}
\right)
\varepsilon_{mn}(\mathbf{x})
+
O(l_{\max}^{4}\nabla^{4}\varepsilon_{mn}) .
\end{equation}

The formulation above is a three-dimensional weakly nonlocal continuum
approximation. Therefore, the one- and two-dimensional finite element
models introduced later should be interpreted as reduced kinematic
specializations of the same three-dimensional theory, rather than as
independently averaged one- or two-dimensional nonlocal models. This
distinction is important because the numerical coefficient in
Eq.~\eqref{eq:nonlocal_strain_result} is determined by the second
moment of a three-dimensional ellipsoidal domain.
\rev{For independently averaged domains the coefficients differ. A 1D segment $[-l,l]$ has $|\mathcal{V}|=2l$ and $\langle\xi^2\rangle=l^2/3$, which gives $1/6$ after the Taylor factor $1/2$. A 2D ellipse $\xi_1^2/l_1^2+\xi_2^2/l_2^2\le1$ has $|\mathcal{V}|=\pi l_1l_2$ with $\langle\xi_i^2\rangle=l_i^2/4$, leading to $1/8$. The 3D ellipsoid gives $|\mathcal{V}|=4\pi/3 l_1l_2l_3$ and $\langle\xi_i^2\rangle=l_i^2/5$, hence $1/10$ (see Appendix A). We retain $1/10$ in 1D and 2D because the physical neighbourhood remains a 3D ellipsoid, while the 1D and 2D models are reduced kinematics (uniaxial and plane-strain) of the same 3D continuum. This choice keeps the models comparable. At $l/L=0.20$ the 1D stiffening ratio is 0.936 and the 2D ratio is 0.94. The 3D ratio is 0.959, similar to the lower-dimensional values. The coefficient $1/10$ has been used consistently across all numerical results in the manuscript (Figs.~1-11). For a truly 1D chain or 2D membrane $1/6$ and $1/8$ would be appropriate. See Appendix A for the full derivation. The energy coefficient $1/20$ yields $1/10$ in the stress after variation by symmetry (see Sec.~2.3).}

%===================================================================
\subsection{Kinematic Relations}
%===================================================================

The displacement field is denoted by $\mathbf{u}(\mathbf{x})$, with
components $u_i$. Under the small-strain assumption, the infinitesimal
strain tensor is
\begin{equation}\label{eq:kinematic}
\varepsilon_{ij}
=
\frac{1}{2}\left(u_{i,j}+u_{j,i}\right),
\end{equation}
where a comma denotes partial differentiation with respect to the
corresponding spatial coordinate. The volumetric strain, or dilatation,
is given by
\begin{equation}\label{eq:dilatation}
\varepsilon_{kk}
=
u_{k,k}
=
\nabla\cdot\mathbf{u}.
\end{equation}

%===================================================================
\subsection{Strain Energy Density and Constitutive Relations}
%===================================================================

Motivated by the anisotropic differential operator obtained from the
ellipsoidal nonlocal averaging procedure, a strain-gradient energy
density is introduced in the form
\begin{equation}\label{eq:energy}
W
=
\frac{1}{2}
C_{ijkl}\varepsilon_{ij}\varepsilon_{kl}
+
\frac{1}{20}
L_{mn}
C_{ijkl}
\varepsilon_{ij,m}\varepsilon_{kl,n},
\qquad
L_{mn}=(AA^T)_{mn}.
\end{equation}
Here $L_{mn}$ is the anisotropic characteristic-length tensor obtained
from the second moment of the ellipsoidal averaging domain. The
coefficient $1/20$ is selected so that the variationally derived stress
correction contains the same second-moment coefficient $1/10$ as the
weakly nonlocal strain expansion, while preserving a positive gradient
contribution to the stored energy. The form of Eq.~\eqref{eq:energy} is
consistent with simplified strain-gradient elasticity models and their
anisotropic extensions
\cite{LazarPo2015,Polizzotto2018}.
\rev{Clarification on Eq.(7) versus Eq.(12) is needed. Eq.~\eqref{eq:nonlocal_strain_result} $\tilde{\varepsilon}=(1+1/10\mathcal{L})\varepsilon$ is a kinematic result. It follows from Taylor expansion of the nonlocal average with a uniform kernel over the ellipsoid. The coefficient $1/10$ is the product of the Taylor factor $1/2$ and the normalized second moment $1/5$. Eq.~\eqref{eq:energy} $W=1/2C\varepsilon\varepsilon+1/20 L_{mn}C\varepsilon_{,m}\varepsilon_{,n}$ is a separate constitutive choice. It is motivated by the operator $L_{mn}$ but is chosen to ensure positive-definiteness with $\mu>0$, $3\lambda+2\mu>0$ and $l_i>0$ giving $L_{mn}$ positive-definite, and to give a well-defined $C^1$ variational formulation. It is not uniquely derived from averaging. Variation gives $\delta W = C_{ijkl}\varepsilon_{kl}\delta\varepsilon_{ij}+1/20 L_{mn}C_{ijkl}(\varepsilon_{ij,m}\delta\varepsilon_{kl,n}+\delta\varepsilon_{ij,m}\varepsilon_{kl,n})$. This becomes $C\varepsilon\delta\varepsilon+1/10 L_{mn}C\varepsilon_{,n}\delta\varepsilon_{,m}$, or $\tau_{ij}\delta\varepsilon_{ij}+q_{ijm}\delta\varepsilon_{ij,m}$ with $q_{ijm}=1/10 L_{mn}C_{ijkl}\varepsilon_{kl,n}$. The form uses symmetry $C_{ijkl}=C_{klij}$ and $L_{mn}=L_{nm}$. Integration by parts gives $\int q_{ijm}\delta\varepsilon_{ij,m}=-\int q_{ijm,m}\delta\varepsilon_{ij}$ plus a boundary term. Hence the effective stress is $\sigma_{ij}=\tau_{ij}-q_{ijm,m}=C\varepsilon-1/10 L_{mn}(C\varepsilon)_{,mn}=(1-1/10\mathcal{L})\sigma^{cl}_{ij}$. Thus $1/20$ in the energy yields $1/10$ in the stress and in the double stress. If $1/10$ were used in the energy, the stress correction would be $1/5$, a double-counting. This separation between kinematic averaging and constitutive regularization clarifies their distinct roles. The local $C_{ijkl}$ is taken isotropic to isolate nonlocal anisotropy coming from the ellipsoidal geometry. Physically this corresponds to materials where the matrix is isotropic but the interaction range is direction-dependent due to elongated features such as fibres, drawn grains and lattice struts. $L_{mn}$ is analogous to the fabric tensor in homogenization \cite{LazarPo2015,Polizzotto2018,Auffray2019}. For materials with additional local anisotropy both sources should be combined as $W=1/2 C^{aniso}\varepsilon\varepsilon+1/20 L_{mn}C^{aniso}\varepsilon_{,m}\varepsilon_{,n}$.}

For the isotropic elastic solid considered here, the fourth-order
elasticity tensor is
\begin{equation}\label{eq:elasticity_tensor}
C_{ijkl}
=
\lambda\delta_{ij}\delta_{kl}
+
\mu\left(\delta_{ik}\delta_{jl}+\delta_{il}\delta_{jk}\right),
\end{equation}
where $\lambda$ and $\mu$ are the Lamé constants. The classical elastic
stress corresponding to the local part of the energy is
\begin{equation}\label{eq:classical_stress}
\sigma_{ij}^{\mathrm{cl}}
=
C_{ijkl}\varepsilon_{kl}
=
\lambda\delta_{ij}\varepsilon_{kk}
+
2\mu\varepsilon_{ij}.
\end{equation}
For $\mu>0$, $3\lambda+2\mu>0$, and positive characteristic lengths
$l_i>0$, the tensor $L_{mn}$ is positive definite and the gradient
contribution in Eq.~\eqref{eq:energy} is non-negative. Thus, the adopted
energy density is positive definite for non-rigid admissible strain
fields.

The total potential energy is written as
\begin{equation}\label{eq:total_potential}
\Pi
=
\int_{\Omega}W\,d\Omega
-
\int_{\Omega}P_i u_i\,d\Omega
-
\int_{\Gamma_t}\bar{t}_i u_i\,d\Gamma ,
\end{equation}
where $P_i$ is the body-force density and $\bar{t}_i$ is the prescribed
classical traction on $\Gamma_t$. Possible higher-order boundary terms
associated with the strain-gradient contribution are not prescribed in
Eq.~\eqref{eq:total_potential}; the corresponding natural boundary
conditions arise from the variational statement and are discussed
separately in the boundary-condition section.

The energetic stress and double stress are defined by
\begin{equation}\label{eq:energetic_stress}
\tau_{ij}
=
\frac{\partial W}{\partial\varepsilon_{ij}}
=
C_{ijkl}\varepsilon_{kl},
\end{equation}
and
\begin{equation}\label{eq:double_stress}
q_{ijm}
=
\frac{\partial W}{\partial\varepsilon_{ij,m}}
=
\frac{1}{10}
L_{mn}
C_{ijkl}\varepsilon_{kl,n}.
\end{equation}
The factor $1/10$ in Eq.~\eqref{eq:double_stress} results from the
symmetry of $C_{ijkl}$ and $L_{mn}$ in the variation of the quadratic
gradient term. For a homogeneous material with constant $C_{ijkl}$ and
constant characteristic-length tensor $L_{mn}$, the effective stress
entering the strong equilibrium equation is
\begin{align}\label{eq:stress_derivation}
\sigma_{ij}
&=
\tau_{ij}
-
q_{ijm,m} \nonumber\\
&=
C_{ijkl}\varepsilon_{kl}
-
\frac{1}{10}
L_{mn}
\frac{\partial^2}{\partial x_m\partial x_n}
\left(C_{ijkl}\varepsilon_{kl}\right) \nonumber\\
&=
\left(1-\frac{1}{10}\mathcal{L}_{\ell}\right)
\sigma_{ij}^{\mathrm{cl}},
\end{align}
where the anisotropic operator $\mathcal{L}_{\ell}$ is defined in
Eq.~\eqref{eq:Lop}. Equivalently,
\begin{equation}\label{eq:stress_explicit}
\sigma_{ij}
=
\left[
1
-
\frac{1}{10}
\left(
l_1^2\frac{\partial^2}{\partial x_1^2}
+
l_2^2\frac{\partial^2}{\partial x_2^2}
+
l_3^2\frac{\partial^2}{\partial x_3^2}
\right)
\right]
\left(
\lambda\delta_{ij}\varepsilon_{kk}
+
2\mu\varepsilon_{ij}
\right).
\end{equation}

The signs in the nonlocal strain approximation and in the effective
stress relation have different origins. The averaged strain obtained from
Taylor expansion satisfies
\begin{equation}\label{eq:strain_sign_remark}
\tilde{\varepsilon}_{ij}
=
\left(1+\frac{1}{10}\mathcal{L}_{\ell}\right)\varepsilon_{ij}
+
O(l_{\max}^{4}\nabla^{4}\varepsilon_{ij}),
\end{equation}
whereas the stress relation in Eq.~\eqref{eq:stress_derivation} follows
from the variation of the positive-definite strain-gradient energy. For a
plane wave $\exp(i\mathbf{k}\cdot\mathbf{x})$,
\[
\mathcal{L}_{\ell}
\rightarrow
-
L_{mn}k_m k_n,
\]
so that
\[
1-\frac{1}{10}\mathcal{L}_{\ell}
\rightarrow
1+\frac{1}{10}L_{mn}k_m k_n .
\]
Hence, the negative sign in the physical-space stress operator produces
a positive stiffness correction in the wavenumber domain.

% ===================================================================
%===================================================================
\subsection{Governing Equation}
%===================================================================

The equation of motion follows from the local balance of linear momentum
written in terms of the effective stress defined in
Eq.~\eqref{eq:stress_derivation}:
\begin{equation}\label{eq:momentum}
\sigma_{ij,j}+P_i
=
\rho\,u_{i,tt},
\end{equation}
where $P_i$ is the body force per unit volume and $\rho$ is the mass
density. For a homogeneous material with constant Lamé parameters and
constant characteristic-length tensor $L_{mn}$, the operator
$\mathcal{L}_{\ell}$ has constant coefficients and commutes with spatial
differentiation. Substituting Eq.~\eqref{eq:stress_explicit} into
Eq.~\eqref{eq:momentum} therefore gives
\begin{equation}\label{eq:governing}
\left(1-\frac{1}{10}\mathcal{L}_{\ell}\right)
\left[
(\lambda+\mu)u_{j,ij}
+
\mu u_{i,jj}
\right]
+
P_i
=
\rho\,u_{i,tt}.
\end{equation}
The term in square brackets is the classical Navier operator,
\begin{equation}\label{eq:classical_navier}
\sigma_{ij,j}^{\mathrm{cl}}
=
(\lambda+\mu)u_{j,ij}
+
\mu u_{i,jj}.
\end{equation}
Equation~\eqref{eq:governing} is therefore a fourth-order anisotropic
strain-gradient extension of the classical elastodynamic equation. Such
higher-order governing equations require additional boundary conditions
beyond the classical displacement or traction conditions
\cite{Mindlin1964,Mindlin1965,AskesAifantis2011}. In the static case,
the inertia term on the right-hand side is omitted.

For reference, consider the isotropic spherical case
$l_1=l_2=l_3=l$, for which
$\mathcal{L}_{\ell}=l^2\nabla^2$. In the absence of body forces, let
\begin{equation}\label{eq:plane_wave}
\mathbf{u}
=
\mathbf{U}\exp\left[i(\mathbf{k}\cdot\mathbf{x}-\omega t)\right],
\end{equation}
where $\mathbf{k}$ is the wave vector, $k=|\mathbf{k}|$, and $\omega$ is
the circular frequency. Substitution into Eq.~\eqref{eq:governing}
yields the longitudinal and transverse dispersion relations
\begin{align}
\text{P-wave:}\qquad
\omega^2
&=
c_P^2 k^2
\left(
1+\frac{l^2 k^2}{10}
\right),
\label{eq:dispersion_P}
\\
\text{S-wave:}\qquad
\omega^2
&=
c_S^2 k^2
\left(
1+\frac{l^2 k^2}{10}
\right),
\label{eq:dispersion_S}
\end{align}
with
\begin{equation}\label{eq:wave_speeds}
c_P
=
\sqrt{\frac{\lambda+2\mu}{\rho}},
\qquad
c_S
=
\sqrt{\frac{\mu}{\rho}}.
\end{equation}
The positive correction factor in
Eqs.~\eqref{eq:dispersion_P}--\eqref{eq:dispersion_S} indicates
short-wavelength stiffening within the range of validity of the
second-order weakly nonlocal approximation, i.e. for $lk\ll1$
\cite{MetrikineAskes2002,AskesAifantis2011}.
\rev{Dispersion relations are included as an analytical stability check. Positive correction for all $k$ confirms positive-definiteness and stability with no imaginary frequencies. In physical space the negative sign $(1- l^2/10\nabla^2)$ becomes positive $1+ l^2k^2/10$ in the wavenumber domain, which explains the sign remark. The result shows short-wavelength stiffening. It is consistent with the static ratio $u_{SG}/u_{cl}\approx1- l/(L\sqrt{10})$ and the dynamic ratio $f_{SG}/f_{cl}>1$. It also indicates validity $lk\ll1$, supporting the $l/L\le0.10$ discussion. No transient wave-propagation simulations are presented; the focus remains static and free vibration. If preferred, this subsection can be moved to Appendix B. For a micro-inertia extension the kinetic energy is $T_g=1/2\int\rho l_k^2\dot{u}_{i,j}\dot{u}_{i,j}d\Omega$. This gives $\omega^2=c^2k^2(1+l_s^2k^2/10)/(1+l_k^2k^2)$ with $c=c_P,c_S$, verified by $\nabla^2\rightarrow -k^2$. For $l_k=0$ the previous result is recovered. For $l_kk\ll1$ the effect is $<5\%$ for the first six modes with $l/L\le0.05$, so it is omitted for low-frequency focus.}

%===================================================================
\section{Particular Cases}
%===================================================================

\subsubsection*{Case 1: Spheroidal neighbourhood}

When $l_1=l_2=l$ and $l_3$ is distinct, the ellipsoidal neighbourhood
has axial symmetry about the $x_3$-axis. The anisotropic operator reduces
to
\begin{equation}\label{eq:spheroidal_operator}
\mathcal{L}_{\ell}
=
l^2\frac{\partial^2}{\partial x_1^2}
+
l^2\frac{\partial^2}{\partial x_2^2}
+
l_3^2\frac{\partial^2}{\partial x_3^2}.
\end{equation}
Accordingly, the effective stress becomes
\begin{equation}\label{eq:spheroidal}
\sigma_{ij}
=
\left[
1
-
\frac{1}{10}
\left(
l^2\frac{\partial^2}{\partial x_1^2}
+
l^2\frac{\partial^2}{\partial x_2^2}
+
l_3^2\frac{\partial^2}{\partial x_3^2}
\right)
\right]
\left(
\lambda\delta_{ij}\varepsilon_{kk}
+
2\mu\varepsilon_{ij}
\right).
\end{equation}
The case $l_3>l$ corresponds to a prolate spheroidal neighbourhood,
whereas $l_3<l$ corresponds to an oblate spheroidal neighbourhood.

\subsubsection*{Case 2: Spherical neighbourhood}

When $l_1=l_2=l_3=l$, the averaging neighbourhood is spherical and
\begin{equation}\label{eq:spherical_operator}
\mathcal{L}_{\ell}
=
l^2\nabla^2.
\end{equation}
The constitutive relation then reduces to
\begin{equation}\label{eq:spherical}
\sigma_{ij}
=
\left(
1-\frac{l^2}{10}\nabla^2
\right)
\left(
\lambda\delta_{ij}\varepsilon_{kk}
+
2\mu\varepsilon_{ij}
\right).
\end{equation}
\rev{Comparison with integral nonlocal, homogenized lattice and alternative gradient formulations is as follows. The integral model is $\tilde{\sigma}_{ij}(\mathbf{x})=\int K(|\mathbf{x}-\mathbf{x}'|)\sigma^{cl}_{ij}(\mathbf{x}')d\Omega'$ with a Gaussian kernel. Taylor expansion gives $\tilde{\sigma}=\sigma^{cl}+l^2/10\nabla^2\sigma^{cl}+O(l^4)$ for a spherical uniform kernel over the ellipsoid. Hence the differential form $(1- l^2/10\nabla^2)\sigma=\sigma^{cl}$ is a second-order approximation with error $O(l^4\nabla^4\sigma)$ (Eringen 1983, Askes \& Aifantis 2011). An anisotropic ellipsoidal kernel gives an anisotropic Laplacian $L_{mn}\partial^2$. Second-order homogenization of lattices \cite{Molavitabrizi2023} shows that the gradient continuum reproduces size effects missed by the classical model, with $L_{mn}\propto\langle\xi_m\xi_n\rangle_{RVE}$. A full comparison would require an RVE of an octet-truss and similar architectures, computing $C_{ijkl}$ and $L_{mn}$ via Hill-Mandel. This is noted as future work requiring a separate solver with real RVE simulations. Compared with Mindlin Form II with 5 lengths and Form III, our simplified Aifantis-type model uses 3 lengths $l_i$ from geometry. It remains positive-definite and is computationally tractable with $C^1$ FEM, whereas full Mindlin requires mixed FEM. The trade-off is less generality but more robustness. The present differential model is thus a weak nonlocal approximation to the integral model, valid for $l_{\max}\ll L_{\varepsilon}$ with error $O(l^4)$.}
This is the isotropic Aifantis-type strain-gradient form
\cite{Aifantis1992,AltanAifantis1997,AskesAifantis2011} with the
effective gradient parameter
\begin{equation}\label{eq:effective_length}
\ell_{\mathrm{eff}}^2
=
\frac{l^2}{10}.
\end{equation}

% ===================================================================
%===================================================================
%===================================================================
\section{Boundary Conditions}
\label{sec:boundary_conditions}
%===================================================================

The governing equation~\eqref{eq:governing} contains fourth-order
spatial derivatives because the anisotropic operator
$\mathcal{L}_{\ell}$ acts on the classical Navier operator. Therefore,
in addition to the classical displacement or traction conditions, one
higher-order boundary condition is required. Let $\boldsymbol{\nu}$ be
the outward unit normal to $\partial\Omega$, and let
$\Gamma_u$ and $\Gamma_t$ denote the displacement and traction parts of
the classical boundary, respectively.

%===================================================================
\subsection{Classical Boundary Conditions}
%===================================================================

On $\Gamma_u$, the displacement is prescribed as
\begin{equation}\label{eq:essential_BC}
u_i=\bar{u}_i
\qquad
\text{on } \Gamma_u .
\end{equation}
On $\Gamma_t$, the effective traction is taken as
\begin{equation}\label{eq:traction_BC}
\sigma_{ij}\nu_j=\bar{t}_i
\qquad
\text{on } \Gamma_t ,
\end{equation}
where $\sigma_{ij}$ is the effective stress defined in
Eq.~\eqref{eq:stress_derivation}. For a traction-free surface,
\begin{equation}\label{eq:free_traction}
\sigma_{ij}\nu_j=0 .
\end{equation}

%===================================================================
\subsection{Higher-Order Boundary Condition}
%===================================================================

For the homogeneous material considered here, the strain-gradient part of
the energy introduces the higher-order stress
\begin{equation}\label{eq:higher_order_stress}
q_{ijm}
=
\frac{\partial W}{\partial \varepsilon_{ij,m}}
=
\frac{1}{10}
L_{mn}C_{ijkl}\varepsilon_{kl,n}
=
\frac{1}{10}
L_{mn}\sigma_{ij,n}^{\mathrm{cl}},
\qquad
L_{mn}=(AA^T)_{mn}.
\end{equation}
The tensor $q_{ijm}$ is work-conjugate to the strain gradient
$\varepsilon_{ij,m}$. The natural higher-order condition is written in
terms of the normal double traction as
\begin{equation}\label{eq:double_traction_BC}
q_{ijm}\nu_m\nu_j=0
\qquad
\text{on } \Gamma_q ,
\end{equation}
where $\Gamma_q\subseteq\partial\Omega$ is the portion of the boundary
on which the higher-order natural condition is imposed. On the
complementary higher-order essential boundary $\Gamma_g$, one may
prescribe the normal derivative of displacement,
\begin{equation}\label{eq:higher_order_essential_BC}
u_{i,\nu}
=
\nabla u_i\cdot\boldsymbol{\nu}
=
\bar{g}_i
\qquad
\text{on } \Gamma_g .
\end{equation}
In the present computations, the natural higher-order condition is used
on free boundaries unless stated otherwise.

Using Eq.~\eqref{eq:higher_order_stress}, Eq.~\eqref{eq:double_traction_BC}
becomes
\begin{equation}\label{eq:BC2}
L_{mn}
\frac{\partial\sigma_{ij}^{\mathrm{cl}}}{\partial x_n}
\nu_m\nu_j
=
0
\qquad
\text{on } \Gamma_q .
\end{equation}
The factor $1/10$ has been omitted since the boundary condition is
homogeneous and multiplication by a nonzero constant does not alter the
resulting constraint. This condition represents the reduced higher-order natural boundary
condition adopted in the present formulation. In accordance with the
underlying Aifantis-type strain-gradient framework 
\cite{Aifantis1992,AltanAifantis1997,AskesAifantis2011}, higher-order
boundary effects are expressed through normal derivatives and normal
double tractions, which are employed consistently throughout the
subsequent finite element formulation.
\rev{Remark on essential versus natural higher-order boundary conditions: Classical conditions are $u_i=\bar{u}_i$ on $\Gamma_u$ and $\sigma_{ij}\nu_j=\bar{t}_i$ on $\Gamma_t$ with $\sigma_{ij}=(1-1/10\mathcal{L})\sigma^{cl}_{ij}$. Higher-order conditions are of two types. (i) Essential: $u_{i,\nu}=\nabla u_i\cdot\boldsymbol{\nu}=\bar{g}_i$ on $\Gamma_g$, enforced through Hermite derivative DOFs $(u_{i,x},u_{i,y},u_{i,z})$. For a clamped-slope bar $\theta(0)=0$ is imposed by setting the nodal derivative DOF to zero. (ii) Natural: $q_{ijm}\nu_m\nu_j=0$ on $\Gamma_q$ with $q_{ijm}=1/10 L_{mn}\sigma^{cl}_{ij,n}$. This follows from the boundary term $\int_{\partial\Omega} q_{ijm}\nu_m\delta\varepsilon_{ij} d\Gamma$ after integration by parts. On free boundaries where derivative DOFs are left free, the weak form enforces this natural condition automatically. For a homogeneous material with constant $L_{mn}$ the operator commutes. The strong form $(1-1/10\mathcal{L})[(\lambda+\mu)u_{j,ij}+\mu u_{i,jj}]+P_i=\rho\ddot{u}_i$ is then fourth-order and needs one extra boundary condition per side, which is provided. For the spherical case on a flat surface $x_3=$const with $\boldsymbol{\nu}=(0,0,1)$, the conditions become $\sigma_{i3}=0$ and $\partial\sigma^{cl}_{i3}/\partial x_3=0$ (Eq.26-27). This is a reduced Aifantis-type condition, not the full Mindlin-Toupin form with surface divergence, noted as a limitation. In the FEM implementation essential higher-order conditions are imposed by constraining derivative DOFs in $K_{ff},F_f$ before scaling. Natural conditions are satisfied variationally. In full Mindlin-Toupin second-gradient theory the variational principle yields surface tractions $t_i$, double tractions $r_i=q_{ijm}\nu_m\nu_j$ and line forces $e_i=\llbracket q_{ijm}\nu_m\mu_j\rrbracket$ along edges where the normal $\nu$ is discontinuous, with $\mu$ the edge co-normal and $\llbracket\cdot\rrbracket$ the jump. For rectangular and brick domains with $90^\circ$ corners the simplified Aifantis-type model adopts reduced conditions $\sigma_{ij}\nu_j=\bar{t}_i$ and $q_{ijm}\nu_m\nu_j=0$. Surface divergence and edge terms are neglected, consistent with the phenomenological energy. Hermite derivative DOFs are shared across elements including corners, so the weak form satisfies natural corner conditions variationally. Reaction error $<10^{-11}$ confirms this. For general polyhedral domains explicit edge conditions are required, noted as a limitation. For a rotated $L_{mn}$ the double traction $q_{ijm}\nu_m=1/10 L_{mn}\sigma^{cl}_{ij,n}\nu_m$ contains off-diagonal contributions. For a clamped face $x=0$ with $\nu=(1,0,0)$, the rotated tensor $L_{rot}=R^T diag(l_x^2,l_y^2,l_z^2)R$ gives $L_{11}=l_x^2c^2+l_y^2s^2$ and $L_{12}=(l_y^2-l_x^2)cs$. Hence $r_{ij}=1/10[L_{11}\sigma^{cl}_{ij,1}+L_{12}\sigma^{cl}_{ij,2}]$. The off-diagonal term couples the transverse gradient $\sigma^{cl}_{ij,2}$ into the double traction even though $\nu$ is along $x$. At $\theta=45^\circ$ the coupling reaches its maximum magnitude with $L_{12}\approx-0.0198$ for $l_x=0.20$ and $l_y=0.02$. This causes energy redistribution, but with the corrected mass $M_{scaled}$ the frequency still decreases monotonically with no minimum at $45^{\circ}$. Tangential components are work-conjugate to $\delta\varepsilon_{ij}$. The natural condition $q_{ijm}\nu_m\nu_j=0$ retains the normal part, while the tangential part would contribute to edge forces in the full theory.}


For a spherical neighbourhood, $L_{mn}=l^2\delta_{mn}$. On a flat
traction-free surface $x_3=\mathrm{const}$ with
$\boldsymbol{\nu}=(0,0,1)$, the classical and higher-order natural
conditions reduce to
\begin{align}
\sigma_{i3}
&=
0,
\qquad i=1,2,3,
\label{eq:flat_free_traction}
\\
\frac{\partial\sigma_{i3}^{\mathrm{cl}}}{\partial x_3}
&=
0,
\qquad i=1,2,3.
\label{eq:flat_double_traction}
\end{align}
The first condition represents the effective traction-free boundary,
while the second condition is the corresponding normal double-traction
condition.



%===================================================================
%===================================================================
\section{Finite Element Formulation}
\label{sec:finite_element_formulation}
%===================================================================

This section presents the displacement-based finite element formulation
corresponding to the strain-gradient model derived above. Since the
governing equation contains fourth-order spatial derivatives, a direct
Galerkin discretization requires $C^1$-continuous basis functions. Cubic
Hermite elements are therefore used in one dimension, Bogner--Fox--Schmit
bicubic Hermite rectangular elements in two dimensions, and
tensor-product tricubic Hermite brick elements in three dimensions
\cite{BognerFoxSchmit1966,PeteraPittman1994,Zervos2008,Papanicolopulos2009}.
The present implementation uses tensor-product rectangular and brick
elements so that Hermite derivative degrees of freedom are shared
consistently across element interfaces.

%===================================================================
\subsection{Weak Form}
\label{subsec:weak_form_fem}
%===================================================================

The strain-gradient energy density is
\begin{equation}
\label{eq:fem_energy_density_general}
W
=
\frac{1}{2}C_{ijkl}\varepsilon_{ij}\varepsilon_{kl}
+
\frac{1}{20}L_{mn}
C_{ijkl}\varepsilon_{ij,m}\varepsilon_{kl,n},
\qquad
L_{mn}=(AA^T)_{mn}.
\end{equation}
For an axis-aligned ellipsoidal neighbourhood,
\begin{equation}
\label{eq:fem_length_tensor}
L_{mn}
=
l_m^2\delta_{mn}
\qquad
\text{no sum on } m .
\end{equation}

Let $\mathbf{u}$ be the displacement field and let $\mathbf{v}$ be an
admissible virtual displacement field vanishing on the essential
boundary. The weak form follows from the first variation of the total
potential energy, supplemented by the inertial virtual work. The problem
is to find $\mathbf{u}\in\mathcal{U}$ such that, for all
$\mathbf{v}\in\mathcal{V}_0$,
\begin{align}
\label{eq:fem_weak_form_general}
&
\int_{\Omega}
\varepsilon_{ij}(\mathbf{v})C_{ijkl}\varepsilon_{kl}(\mathbf{u})
\,d\Omega
+
\frac{1}{10}
\int_{\Omega}
L_{mn}
\varepsilon_{ij,m}(\mathbf{v})
C_{ijkl}
\varepsilon_{kl,n}(\mathbf{u})
\,d\Omega
\nonumber\\
&\qquad
+
\int_{\Omega}
\rho\,v_i\ddot{u}_i\,d\Omega
=
\int_{\Omega}v_iP_i\,d\Omega
+
\int_{\Gamma_t}v_i\bar{t}_i\,d\Gamma .
\end{align}
For static problems, the inertial term is omitted. The higher-order
natural boundary terms are those associated with the boundary conditions
introduced in Section~\ref{sec:boundary_conditions}.

After spatial discretization, the semi-discrete finite element equations
take the standard form
\begin{equation}
\label{eq:fem_semidiscrete}
\mathbf{M}\ddot{\mathbf{d}}
+
\mathbf{K}\mathbf{d}
=
\mathbf{F},
\end{equation}
where
\begin{equation}
\label{eq:fem_total_stiffness}
\mathbf{K}
=
\mathbf{K}_{\mathrm{cl}}
+
\mathbf{K}_{\mathrm{g}} .
\end{equation}
The classical stiffness matrix is
\begin{equation}
\label{eq:fem_classical_stiffness}
\mathbf{K}_{\mathrm{cl}}
=
\int_{\Omega}\mathbf{B}^{T}\mathbf{C}\mathbf{B}\,d\Omega,
\end{equation}
where $\mathbf{B}$ is the classical strain-displacement matrix. The
gradient stiffness matrix is
\begin{equation}
\label{eq:fem_gradient_stiffness_general}
\mathbf{K}_{\mathrm{g}}
=
\frac{1}{10}
\int_{\Omega}
\sum_{\alpha=1}^{d}
\sum_{\beta=1}^{d}
L_{\alpha\beta}
\left(\frac{\partial\mathbf{B}}{\partial x_{\alpha}}\right)^{T}
\mathbf{C}
\left(\frac{\partial\mathbf{B}}{\partial x_{\beta}}\right)
\,d\Omega ,
\end{equation}
where $d$ is the spatial dimension of the reduced kinematic model.
For an axis-aligned ellipsoidal neighbourhood,
Eq.~\eqref{eq:fem_gradient_stiffness_general} reduces to
\begin{equation}
\label{eq:fem_gradient_stiffness_aligned}
\mathbf{K}_{\mathrm{g}}
=
\frac{1}{10}
\int_{\Omega}
\sum_{\alpha=1}^{d}
l_{\alpha}^{2}
\left(\frac{\partial\mathbf{B}}{\partial x_{\alpha}}\right)^{T}
\mathbf{C}
\left(\frac{\partial\mathbf{B}}{\partial x_{\alpha}}\right)
\,d\Omega .
\end{equation}
For the spherical characteristic-length case, $l_1=l_2=l_3=l$, this becomes
\begin{equation}
\label{eq:fem_gradient_stiffness_spherical}
\mathbf{K}_{\mathrm{g}}
=
\frac{l^2}{10}
\int_{\Omega}
\sum_{\alpha=1}^{d}
\left(\frac{\partial\mathbf{B}}{\partial x_{\alpha}}\right)^{T}
\mathbf{C}
\left(\frac{\partial\mathbf{B}}{\partial x_{\alpha}}\right)
\,d\Omega .
\end{equation}

The consistent mass matrix is
\begin{equation}
\label{eq:fem_mass_general}
\mathbf{M}
=
\int_{\Omega}
\rho\,\mathbf{N}^{T}\mathbf{N}\,d\Omega ,
\end{equation}
where $\mathbf{N}$ is the displacement interpolation matrix. For
free-vibration analysis, Eq.~\eqref{eq:fem_semidiscrete} gives the
generalized eigenvalue problem
\begin{equation}
\label{eq:fem_eigenproblem_general}
\mathbf{K}\boldsymbol{\phi}
=
\omega^2\mathbf{M}\boldsymbol{\phi}.
\end{equation}

%===================================================================
\subsection{One-Dimensional Hermite Element}
\label{subsec:one_dimensional_hermite_element}
%===================================================================

For the one-dimensional bar, the displacement field is $u(x)$ and
\begin{equation}
\label{eq:1d_strain_gradient}
\varepsilon=u',
\qquad
\varepsilon_{,x}=u'' .
\end{equation}
The strain energy reduces to
\begin{equation}
\label{eq:1d_fem_energy}
U
=
\frac{1}{2}
\int_0^L
EA_c\left[(u')^2+a^2(u'')^2\right]dx,
\qquad
a^2=\frac{l^2}{10},
\end{equation}
where $A_c$ is the cross-sectional area.

A two-node cubic Hermite element is used. Each node contains the
displacement and its first derivative,
\begin{equation}
\label{eq:1d_dofs}
\mathbf{d}_i
=
\begin{bmatrix}
u_i & \theta_i
\end{bmatrix}^{T},
\qquad
\theta_i
=
\left.\frac{du}{dx}\right|_{x=x_i}.
\end{equation}
For an element of length $L_e$,
\begin{equation}
\label{eq:1d_hermite_interpolation}
u^e(x)
=
\mathbf{N}\mathbf{d}^e,
\qquad
\mathbf{d}^e
=
\begin{bmatrix}
u_1 & \theta_1 & u_2 & \theta_2
\end{bmatrix}^{T}.
\end{equation}
Using the local coordinate $s=x/L_e$, $0\leq s\leq1$, the Hermite shape
functions are
\begin{align}
N_1 &= 1-3s^2+2s^3, \label{eq:1d_H_N1}\\
N_2 &= L_e(s-2s^2+s^3), \label{eq:1d_H_N2}\\
N_3 &= 3s^2-2s^3, \label{eq:1d_H_N3}\\
N_4 &= L_e(-s^2+s^3). \label{eq:1d_H_N4}
\end{align}
Defining
\begin{equation}
\label{eq:1d_B_matrices}
\mathbf{B}_1
=
\frac{d\mathbf{N}}{dx},
\qquad
\mathbf{B}_2
=
\frac{d^2\mathbf{N}}{dx^2},
\end{equation}
the element stiffness matrix is
\begin{equation}
\label{eq:1d_element_stiffness}
\mathbf{K}^{e}
=
\int_{x_e}
EA_c
\left(
\mathbf{B}_1^T\mathbf{B}_1
+
a^2\mathbf{B}_2^T\mathbf{B}_2
\right)dx .
\end{equation}
The consistent mass matrix is
\begin{equation}
\label{eq:1d_element_mass}
\mathbf{M}^{e}
=
\int_{x_e}
\rho A_c\,\mathbf{N}^T\mathbf{N}\,dx .
\end{equation}

%===================================================================
\subsection{Two-Dimensional Bogner--Fox--Schmit Element}
\label{subsec:two_dimensional_bfs_element}
%===================================================================

For the two-dimensional plane-strain formulation, the displacement field
is
\begin{equation}
\label{eq:2d_displacement}
\mathbf{u}
=
\begin{bmatrix}
u & v
\end{bmatrix}^{T},
\end{equation}
and the engineering strain vector is
\begin{equation}
\label{eq:2d_fem_strain_vector}
\boldsymbol{\varepsilon}
=
\begin{bmatrix}
\varepsilon_{xx} \\
\varepsilon_{yy} \\
\gamma_{xy}
\end{bmatrix}
=
\begin{bmatrix}
u_{,x} \\
v_{,y} \\
u_{,y}+v_{,x}
\end{bmatrix}.
\end{equation}
The matrix $\mathbf{C}$ in this subsection denotes the corresponding
plane-strain elasticity matrix written in Voigt form consistent with the
engineering shear strain convention.

A Bogner--Fox--Schmit bicubic Hermite rectangular element is used for
each displacement component. For a scalar field $w$, the nodal degrees of
freedom are
\begin{equation}
\label{eq:2d_scalar_bfs_dofs}
w,
\qquad
w_{,x},
\qquad
w_{,y},
\qquad
w_{,xy}.
\end{equation}
Thus, each node has eight degrees of freedom in two-dimensional
elasticity:
\begin{equation}
\label{eq:2d_bfs_dofs_fem}
\mathbf{d}_i
=
\begin{bmatrix}
u_i & u_{i,x} & u_{i,y} & u_{i,xy} &
v_i & v_{i,x} & v_{i,y} & v_{i,xy}
\end{bmatrix}^{T}.
\end{equation}

For the spherical characteristic-length case, the element stiffness matrix is
\begin{equation}
\label{eq:2d_bfs_stiffness_isotropic}
\mathbf{K}^{e}
=
\int_{\Omega_e}
\left[
\mathbf{B}^{T}\mathbf{C}\mathbf{B}
+
a^2
\left(
\mathbf{B}_{,x}^{T}\mathbf{C}\mathbf{B}_{,x}
+
\mathbf{B}_{,y}^{T}\mathbf{C}\mathbf{B}_{,y}
\right)
\right]d\Omega,
\qquad
a^2=\frac{l^2}{10}.
\end{equation}
For an axis-aligned ellipsoidal model, the stiffness matrix becomes
\begin{equation}
\label{eq:2d_bfs_stiffness_anisotropic}
\mathbf{K}^{e}
=
\int_{\Omega_e}
\left[
\mathbf{B}^{T}\mathbf{C}\mathbf{B}
+
\frac{l_1^2}{10}
\mathbf{B}_{,x}^{T}\mathbf{C}\mathbf{B}_{,x}
+
\frac{l_2^2}{10}
\mathbf{B}_{,y}^{T}\mathbf{C}\mathbf{B}_{,y}
\right]d\Omega .
\end{equation}
If a rotated characteristic-length tensor is used, the cross term
$L_{12}\mathbf{B}_{,x}^{T}\mathbf{C}\mathbf{B}_{,y}$ and its symmetric
counterpart follow directly from Eq.~\eqref{eq:fem_gradient_stiffness_general}.

The consistent mass matrix is
\begin{equation}
\label{eq:2d_bfs_mass}
\mathbf{M}^{e}
=
\int_{\Omega_e}
\rho\,\mathbf{N}_v^T\mathbf{N}_v\,d\Omega ,
\end{equation}
where $\mathbf{N}_v$ is the vector displacement interpolation matrix.

%===================================================================
\subsection{Three-Dimensional Tricubic Hermite Brick Element}
\label{subsec:three_dimensional_hermite_element}
%===================================================================

For the three-dimensional formulation, the displacement vector is
\begin{equation}
\label{eq:3d_displacement}
\mathbf{u}
=
\begin{bmatrix}
u & v & w
\end{bmatrix}^{T},
\end{equation}
and the engineering strain vector is
\begin{equation}
\label{eq:3d_fem_strain_vector}
\boldsymbol{\varepsilon}
=
\begin{bmatrix}
\varepsilon_{xx} \\
\varepsilon_{yy} \\
\varepsilon_{zz} \\
\gamma_{yz} \\
\gamma_{xz} \\
\gamma_{xy}
\end{bmatrix}
=
\begin{bmatrix}
u_{,x} \\
v_{,y} \\
w_{,z} \\
v_{,z}+w_{,y} \\
u_{,z}+w_{,x} \\
u_{,y}+v_{,x}
\end{bmatrix}.
\end{equation}
Here also, $\mathbf{C}$ is understood in Voigt form consistent with the
engineering shear strain convention.

A tensor-product tricubic Hermite brick element is used. For each scalar
field $q$, the nodal degrees of freedom are
\begin{equation}
\label{eq:3d_scalar_hermite_dofs}
q,
\quad
q_{,x},
\quad
q_{,y},
\quad
q_{,z},
\quad
q_{,xy},
\quad
q_{,xz},
\quad
q_{,yz},
\quad
q_{,xyz}.
\end{equation}
Thus, for the three displacement components, each node contains
$24$ degrees of freedom:
\begin{equation}
\label{eq:3d_hermite_dofs_fem}
\begin{aligned}
\mathbf{d}_i=[&
u_i,u_{i,x},u_{i,y},u_{i,z},u_{i,xy},u_{i,xz},u_{i,yz},u_{i,xyz},\\
&
v_i,v_{i,x},v_{i,y},v_{i,z},v_{i,xy},v_{i,xz},v_{i,yz},v_{i,xyz},\\
&
w_i,w_{i,x},w_{i,y},w_{i,z},w_{i,xy},w_{i,xz},w_{i,yz},w_{i,xyz}
]^T .
\end{aligned}
\end{equation}

For the spherical characteristic-length case, the element stiffness matrix is
\begin{equation}
\label{eq:3d_hermite_stiffness_isotropic}
\mathbf{K}^{e}
=
\int_{\Omega_e}
\left[
\mathbf{B}^{T}\mathbf{C}\mathbf{B}
+
a^2
\left(
\mathbf{B}_{,x}^{T}\mathbf{C}\mathbf{B}_{,x}
+
\mathbf{B}_{,y}^{T}\mathbf{C}\mathbf{B}_{,y}
+
\mathbf{B}_{,z}^{T}\mathbf{C}\mathbf{B}_{,z}
\right)
\right]d\Omega,
\qquad
a^2=\frac{l^2}{10}.
\end{equation}
For the axis-aligned ellipsoidal model with distinct characteristic
lengths, the stiffness matrix is
\begin{equation}
\label{eq:3d_hermite_stiffness_anisotropic}
\mathbf{K}^{e}
=
\int_{\Omega_e}
\left[
\mathbf{B}^{T}\mathbf{C}\mathbf{B}
+
\frac{l_1^2}{10}
\mathbf{B}_{,x}^{T}\mathbf{C}\mathbf{B}_{,x}
+
\frac{l_2^2}{10}
\mathbf{B}_{,y}^{T}\mathbf{C}\mathbf{B}_{,y}
+
\frac{l_3^2}{10}
\mathbf{B}_{,z}^{T}\mathbf{C}\mathbf{B}_{,z}
\right]d\Omega .
\end{equation}
For a rotated ellipsoidal neighbourhood, the general tensor form in
Eq.~\eqref{eq:fem_gradient_stiffness_general} is used.

The consistent mass matrix is
\begin{equation}
\label{eq:3d_hermite_mass}
\mathbf{M}^{e}
=
\int_{\Omega_e}
\rho\,\mathbf{N}_v^T\mathbf{N}_v\,d\Omega .
\end{equation}

%===================================================================
\subsection{Implementation Remarks}
\label{subsec:fem_implementation_remarks}
%===================================================================

Hermite-type elements contain displacement derivatives as nodal degrees
of freedom. The corresponding stiffness and mass matrices may therefore
contain entries with different physical dimensions and scales. To improve
conditioning, a symmetric diagonal scaling is applied to the free-degree
system. For static analysis,
\begin{equation}
\label{eq:diagonal_scaling}
\tilde{\mathbf{K}}
=
\mathbf{S}\mathbf{K}_{ff}\mathbf{S},
\qquad
\tilde{\mathbf{F}}
=
\mathbf{S}\mathbf{F}_{f},
\qquad
\mathbf{d}_{f}
=
\mathbf{S}\tilde{\mathbf{d}},
\end{equation}
where $\mathbf{K}_{ff}$ and $\mathbf{F}_f$ are the stiffness matrix and
force vector after imposing essential boundary conditions. In the present
implementation, $\mathbf{S}$ is chosen as a diagonal matrix based on the
inverse square root of the positive diagonal entries of $\mathbf{K}_{ff}$.

For free-vibration analysis, the scaled generalized eigenvalue problem is
\begin{equation}
\label{eq:scaled_eigenproblem}
\tilde{\mathbf{K}}\tilde{\boldsymbol{\phi}}
=
\omega^2
\tilde{\mathbf{M}}\tilde{\boldsymbol{\phi}},
\qquad
\tilde{\mathbf{M}}
=
\mathbf{S}\mathbf{M}_{ff}\mathbf{S},
\end{equation}
and the physical eigenvector is recovered from
\begin{equation}
\label{eq:eigenvector_recovery}
\boldsymbol{\phi}
=
\mathbf{S}\tilde{\boldsymbol{\phi}}.
\end{equation}
The scaling changes only the numerical representation of the algebraic
system; it does not modify the finite element approximation.

The essential boundary conditions in the direct strain-gradient
implementation include classical displacement constraints and, when
prescribed, Hermite derivative constraints. On boundaries where the
corresponding derivative degrees of freedom are left free, the natural
higher-order boundary conditions are enforced through the weak form.
\rev{The present formulation is given first in general 3D form. The weak form is $\int\varepsilon(v)C\varepsilon(u)+1/10\int L_{mn}\varepsilon_{,m}(v)C\varepsilon_{,n}(u)$ with stiffness $K=K_{cl}+K_g$. The gradient part is $K_g=1/10\int\sum L_{\alpha\beta}(\partial B/\partial x_\alpha)^T C (\partial B/\partial x_\beta)$. One-dimensional cubic Hermite, two-dimensional BFS bicubic and three-dimensional tricubic elements are reduced kinematic specializations of the same 3D continuum, corresponding to uniaxial $u(x)$, plane-strain $u(x,y)$ and full $u(x,y,z)$. The 3D coefficient $1/10$ is retained for cross-dimensional consistency, as detailed in Appendix A. For truly lower-dimensional materials $1/6$ for 1D and $1/8$ for 2D would be appropriate. Extension to irregular geometries requires isoparametric $C^1$ Hermite elements. The parent domain $(\xi,\eta,\zeta)\in[-1,1]^3$ is mapped to the physical domain via Hermite blending that preserves $C^1$ continuity. Derivatives are transformed via the Jacobian $J=\partial\mathbf{x}/\partial\boldsymbol{\xi}$ with the Hessian for second derivatives, and $C^1$ continuity of the normal derivative across curved interfaces is enforced via constraints or BFS blending. For general quadrilateral and hexahedral meshes $C^1$ continuity is non-trivial and requires Argyris or Bell elements, or alternatively isogeometric analysis with NURBS basis ($C^{p-1}$) or virtual element formulations \cite{WriggersHudobivnik2023}. The present rectangular formulation serves as a verification basis, with curved boundaries left as future work. For 3D mesh sensitivity, tricubic Hermite elements contain derivative DOFs. The clamped-gradient conditions $u_i=0$ and $u_{i,j}=0$ at $x=0$ are enforced essentially, which allows an exponential boundary layer $\theta(x)\approx F/EA(1-e^{-x/a})$ with $a=\sqrt{l^2/10}=0.0158$ m for $l=0.05$ m. For a $2\times1\times1$ mesh $h_x=0.5$ m $\gg a$, the boundary layer is under-resolved. This causes $5.9\%$ displacement variation and $9.3\%$ frequency variation compared with a $6\times3\times3$ mesh with $h_x=0.166$ m. Real output from \texttt{mesh\_sensitivity\_3D.m} shows displacement increasing from $u=4.5665$ $\mu$m at $2\times1\times1$ to $u=4.8371$ $\mu$m at $6\times3\times3$, while frequency decreases from $f_1=193.00$ Hz to $f_1=174.96$ Hz, with the intermediate $4\times2\times2$ mesh giving $u=4.7876$ $\mu$m and $f_1=178.72$ Hz. Convergence is monotonic from below for displacement and from above for frequency, consistent with a conforming over-stiff formulation. A fourth-order rate is expected only when $h_x \ll a$, which requires more than 20 elements in the x-direction for $l/L_x=0.05$. Gradient energy is integrated with $3\times3\times3$ Gauss points exactly. Diagonal scaling keeps $\kappa<10^6$ and equilibrium error $<10^{-11}$. With respect to micro-inertia, the present model uses classical kinetic energy $T=1/2\int\rho\dot{u}_i\dot{u}_i d\Omega$ and omits $T_g=1/2\int\rho l_k^2\dot{u}_{i,j}\dot{u}_{i,j}d\Omega$ (Mindlin 1964). Including $T_g$ gives an enriched mass $M_g$ and dispersion $\omega^2=c^2k^2(1+l_s^2k^2/10)/(1+l_k^2k^2)$. For weakly nonlocal response $l_kk\ll1$ and low-frequency first six modes with $k_nL\sim(2n-1)\pi/2$, the condition $lk_n<0.3$ holds for $l/L\le0.05$ and the effect is $<5\%$, so it is omitted for clarity.}



\section{Numerical Results and Discussion}
\label{sec:numerical_results}


% =========================================================================
%  5.1  MATERIAL AND SIMULATION PARAMETERS
% =========================================================================
\subsection{Material and Simulation Parameters}
\label{sec:materials}

Unless otherwise specified, all numerical results presented in this
study are obtained for a structural steel characterised by a Young's
modulus of $E = 200\,\mathrm{GPa}$ and a mass density of
$\rho = 7800\,\mathrm{kg\,m^{-3}}$.

For adopted parameters for the one dimensional bar example are $L = 1\,\mathrm{m}$, $A = 1\times10^{-4}\,\mathrm{m^{2}}$, and
$F = 10^{4}\,\mathrm{N}$. Whereas for a two dimensional plane-strain domain example a domain of size $L_{x}\times L_{y} = 1\,\mathrm{m}\times 0.2\,\mathrm{m}$ is considered with the Poisson
ratio is $\nu = 0.30$ and uniform right edge traction 
$t_{x} = 10^{6}\,\mathrm{Pa}$. In order to examine size effects, the ratio $l/L$ is chosen in the scale of $0$ to $0.20$. \rev{Note: $l_i$ are material constants (fabric-tensor-like second moments), while anisotropy ratio $AR=l_x/l_y$ and orientation $\theta$ are processing-induced design manifestations (fibre orientation, print direction, drawing direction), not optimization variables tuned with external fields. See Sec.\ref{sec:calibration}.} A boundary layer reference case with $l/L=0.05$ has been considered unless otherwise noted. These settings help to expose the exponential boundary layer profile  and maintain a global displacement response. 



% For the one-dimensional bar, the geometric and loading parameters are
% $L = 1\,\mathrm{m}$, $A = 1\times10^{-4}\,\mathrm{m^{2}}$, and
% $F = 10^{4}\,\mathrm{N}$.
% For the two-dimensional plane-strain plate, the domain is
% $L_{x}\times L_{y} = 1\,\mathrm{m}\times 0.2\,\mathrm{m}$, the Poisson
% ratio is $\nu = 0.30$, and the uniform right-edge traction is
% $t_{x} = 10^{6}\,\mathrm{Pa}$.
% The microstructural length parameter is varied over the range
% $l/L \in [0,\,0.20]$ in the parametric studies; the boundary-layer
% reference case uses $l/L = 0.05$, corresponding to a gradient length
% $a = \sqrt{l^{2}/10}\approx 1.6\times10^{-2}\,\mathrm{m}$, which
% clearly reveals the exponential boundary-layer structure while keeping
% the gradient contribution modest enough that the global displacement
% remains close to the classical solution.
% These parameters are adopted throughout unless stated otherwise.




The one-dimensional problem is treated first since the clamped slope strain gradient bar \rev{admits a closed-form analytical solution} enabling rigorous verification of the proposed $C^{1}$ Hermite finite
element implementation.  Following the verification for one dimension, the formulation is extended to two and three dimensions via Bogner--Fox--Schmit (BFS) bicubic Hermite element, and to three-dimensional problems using the tensor-product tricubic Hermite element, respectively.  These extensions preserves the same inter-element continuity conditions and the gradient stiffness constructions.




% =========================================================================
%  5.2  ONE-DIMENSIONAL STRAIN-GRADIENT BAR
% =========================================================================
\subsection{One-Dimensional Strain-Gradient Bar}
\label{sec:1D_bar}

% -----------------------------------------------------------------
%  5.2.1  Boundary-Layer Characteristics
% -----------------------------------------------------------------
\subsubsection{Boundary-Layer Characteristics}
\label{sec:1D_boundary_layer}

\begin{figure}[!htbp]
    \centering
    \includegraphics[width=\linewidth]{Fig01_1D_BoundaryLayer.eps}
    \caption{Boundary-layer characteristics of the one-dimensional
    strain-gradient bar under clamped-slope conditions
    ($l/L = 0.05$, $E = 2\times10^{11}$\,Pa,
    $A = 10^{-4}$\,m$^{2}$, $L = 1$\,m,
    $F = 10^{4}$\,N, $N_{\mathrm{el}} = 100$):
    (a)~axial displacement distribution $u(x)$;
    (b)~slope distribution $\theta(x) = \mathrm{d}u/\mathrm{d}x$.
    Displacement remains nearly linear, while slope shows exponential layer near clamped end.}
    \label{fig:1D_BoundaryLayer}
\end{figure}

The displacements and their slopes predicted using the proposed
$C^{1}$ Hermite finite element method for the strain gradient bar under
tip load with essential boundary condition $\theta(0)=0$ applied at the
fixed end are shown in Figure~\ref{fig:1D_BoundaryLayer}.
As can be seen from the two plots, the solution uses two different
length scales: a globally nearly linear displacement that is determined
by the classical stiffness, and an exponentially decaying local
correction to the strain field in a thin layer near the clamped end.






It can be seen from panel~(a) that the displacement is a monotonic function,
starting from $u(0) = 0$ up to the end subjected to loading.
At the scale of the entire domain the curve is visually indistinguishable from the classical one,
$u_{\mathrm{cl}}(x) = Fx/(EA)$, as the axial force is everywhere constant and thus would produce
uniform strain in the case of no additional constraints.
Analytical value of the displacement at the loaded end in the case of a clamped-slope problem
subjected to a tip force is $u_{\mathrm{SG}}(L) = (F/EA)(L - a\tanh(L/a))$, resulting
in an approximate value of $492\,\mu\mathrm{m}$ --- some $1.6\,\%$ less than the classical
one, $FL/(EA) =  500\,\mu\mathrm{m}$.

The effect cannot be observed on the scale of panel~(a),
but it is important in that it constitutes an observable effect of the
microstructural length parameter: it is non-zero for any $l > 0$ and
is missing in any classical analysis, regardless of the fineness of the mesh used.
The observed stiffening is a direct consequence of the higher-order slope constraint $\theta(0)=0$, which forces the strain distribution to depart from uniformity near the fixed end, thereby activating the strain-gradient energy contribution. If natural boundary conditions are imposed (i.e., no constraint on
the slope at either end), then the tip-loaded strain gradient bar
has the classical uniform strain solution as its exact solution.


Slope distribution depicted in panel~(b) illustrates this process
directly.
According to classical theory of elasticity, axial strain
$\varepsilon = \mathrm{d}u/\mathrm{d}x$ under a constant internal
resultant of force is uniform and equals $F/(EA)$. The second-order
governing equation does not have a separate essential boundary
condition for slope and cannot generate any boundary layer.
The fourth-order governing equation with respect to the displacement
allows imposing two essential boundary conditions independently of
the displacement boundary conditions.
The slope boundary condition $\theta(0)=0$ corresponds precisely to
this case: it is possible to impose this condition in the context of
strain gradient theory but has no counterpart in classical theory.


The slope grows from zero value at $x=0$, goes through the sharp monotonic
increase and approaches the classical plateau $F/(EA)=5 \times 10^{-4}$
when $x/L \lesssim 0.10$. 
Then, $\theta(x)$ becomes constant to the precision of the graph
plotting for the rest of the area. 

The length of the boundary layer is determined by the gradient length 
$a = \sqrt{l^{2}/10}$ resulting from the ellipsoidal non-local averaging procedure.

%
\begin{equation}
    \theta(x) \approx \frac{F}{EA}\!\left(1 - e^{-x/a}\right),
    \label{eq:theta_approx}
\end{equation}
%
Given the current parameters, $a/L \approx 1.6 \times 10^{-2}$, which means that the layer covers approximately the first two percent of the total length of the bar. Given the assumption that $L/a \gg 1$, which applies for the current case $(L/a \approx 63)$,
the exact analytical solution reduces to the one which shows that the boundary layer is governed by an exponential
kernel with length scale equal to $a$.

The classical result is obtained in the limiting case of $l \to 0$, as $a \to 0$ implies $e^{-x/a} \to 0$ for all $x > 0$.
The finite element results agree with the analytical slope profile
to $10^{-3}\,\%$ accuracy in the case of the present 100-element
mesh; the relative tip-displacement error reduces from about
$1.2\,\%$ in the case of a coarse 5-element mesh to less than
$10^{-3}\,\%$ in the present 100-element mesh, indicating the
fourth order of algebraic convergence of the $C^{1}$ cubic Hermite
interpolation which is expected and proving that the layer structure
is well resolved with a uniformly-spaced mesh when $h \ll a$.

Together, both panels show that the strain-gradient influences act within the
boundary zones due to higher-order boundary conditions, whereas the
interior behavior follows classic elasticity principles. This separation of
scales is an inherent property of Aifantis-type gradient elastic models.
Such behavior of solutions leads us to the parametric study of the
microstructural parameter.

% -----------------------------------------------------------------
%  5.2.2  Influence of the Microstructural Length Parameter
%         on Static Response
% -----------------------------------------------------------------
\subsubsection{Influence of the Microstructural Length Parameter on
               Static Response}
\label{sec:micro_length_effect}

\begin{figure*}[!htbp]
    \centering
    \includegraphics[width=\linewidth]{Fig02_MicroLengthEffect.eps}
    \caption{Influence of the microstructural length parameter on the
    static response of the one-dimensional strain-gradient bar
    ($E = 2\times10^{11}$\,Pa,
    $A = 10^{-4}$\,m$^{2}$, $L = 1$\,m,
    $F = 10^{4}$\,N):
    (a)~tip displacement $u(L)$;
    (b)~normalised displacement ratio
    $u_{\mathrm{SG}}(L)/u_{\mathrm{cl}}(L)$ versus $l/L$.
    Both quantities decrease monotonically with $l/L$.}
    \label{fig:micro_length_effect}
\end{figure*}

Figure~\ref{fig:micro_length_effect} presents the results of a
parametric study in which the microstructural length parameter $l$ is
varied over the range $l/L \in [0,\,0.20]$ while all other material
and geometric parameters are held constant.
Panel~(a) demonstrates that the tip displacement $u(L)$ decreases
monotonically as $l/L$ increases, departing progressively from the
classical prediction $u_{\mathrm{cl}}(L) = FL/(EA) = 500\,\mu\mathrm{m}$
and reaching approximately $468\,\mu\mathrm{m}$ at $l/L = 0.20$.
This behavior shows that this stiffening is not a change in the classical
elastic modulus, but the activation of a higher-order resistance to
non-uniform strain distributions.
The clamped-slope condition $\theta(0) = 0$ forces the axial strain to
depart from the uniform state near the fixed boundary, and the
strain-gradient energy term
%
\begin{equation*}
    \tfrac{1}{2}EA\,a^{2}\int_{0}^{L}(u'')^{2}\,\mathrm{d}x
\end{equation*}
%
penalises the resulting strain curvature with a contribution that
scales with $a^{2} = l^{2}/10$.
This additional stiffness grows with the microstructural length
parameter $l$, which characterises the spatial extent of the nonlocal
microstructural averaging domain and is a property of the material
rather than of the bar geometry.
Consequently, the effective structural stiffness increases with $l$,
and the corresponding tip displacement falls below the classical
prediction by an amount that grows with $l/L$.

The normalised displacement ratio
$u_{\mathrm{SG}}(L)/u_{\mathrm{cl}}(L)$ shown in panel~(b) provides a
dimensionless measure of the stiffening that is independent of the
applied load $F$ and the cross-sectional stiffness $EA$, depending
instead solely on the dimensionless ratio $l/L$.
The ratio equals unity at $l/L = 0$, confirming that the
strain-gradient model converges to classical elasticity when the
microstructural length vanishes, and it decreases with a near-linear
dependence on $l/L$ throughout the range considered.
This near-linearity follows from the analytical solution
$u_{\mathrm{SG}}(L) = (F/EA)(L - a\tanh(L/a))$: for $L/a \gg 1$
(satisfied for all values shown), $\tanh(L/a) \to 1$ and the ratio
reduces to $1 - l/(L\sqrt{10})$, a linear function of $l/L$ with
slope $-1/\sqrt{10}$.
At $l/L = 0.10$, the normalised ratio is approximately $0.968$,
indicating a tip displacement roughly $3.2\,\%$ below the classical
value --- a difference that is systematic, repeatable, and absent from
any classical finite element analysis regardless of mesh refinement.
The agreement between the $C^{1}$ Hermite finite element solution and
the exact analytical result is excellent in both panels, consistent
with the fourth-order convergence established in
Section~\ref{sec:1D_boundary_layer}.

These results confirm that the proposed one-dimensional discretisation
is both sufficiently accurate and mesh-converged.
Having established the size-dependent static response, we now examine
the free-vibration characteristics of the strain-gradient bar.


% -----------------------------------------------------------------
%  5.2.3  Influence of the Microstructural Length Parameter
%         on Free-Vibration Characteristics
% -----------------------------------------------------------------
\subsubsection{Influence of the Microstructural Length Parameter on
               Free-Vibration Characteristics}
\label{sec:eigenfrequency}

\begin{figure*}[!htbp]
    \centering
    \includegraphics[width=\linewidth]{Fig03_EigenfrequencyResponse.eps}
    \caption{Free-vibration characteristics of the one-dimensional
    strain-gradient bar
    ($E = 2\times10^{11}\,\mathrm{Pa}$,
    $\rho = 7800\,\mathrm{kg\,m^{-3}}$,
    $A = 10^{-4}\,\mathrm{m}^{2}$, $L = 1\,\mathrm{m}$,
    $N_{\mathrm{el}} = 120$):
    (a)~first six natural frequencies versus $l/L$;
    (b)~normalised frequency ratio
    $f_{\mathrm{SG}}/f_{\mathrm{classical}}$ versus $l/L$.
    Every natural frequency increases monotonically with $l/L$, with
    higher modes exhibiting progressively greater strain-gradient
    stiffening.}
    \label{fig:eigenfrequency_response}
\end{figure*}

Figure~\ref{fig:eigenfrequency_response} presents the free-vibration
response of the strain-gradient bar as a function of the normalised
microstructural length $l/L$ over the range $[0,\,0.20]$.
Panel~(a) shows that every natural frequency increases monotonically
with $l/L$, so the strain-gradient formulation predicts a stiffer
dynamic response than classical elasticity for all modes and all
non-zero values of the microstructural length.
The increase is markedly mode-dependent across the spectrum. Notably Mode~1 increases only
modestly from nearly $1.3\,\mathrm{kHz}$ to
$1.4\,\mathrm{kHz}$ --- whereas Mode~6 experiences a substantial 
elevation, growing from roughly $14\,\mathrm{kHz}$ to more than
$21\,\mathrm{kHz}$ over the same parameter range.
The growth patterns become increasingly nonlinear as $l/L$ increases, with the
modal patterns bending more sharply at larger values of the microstructural length.
This concave-upward behaviour of the curves is consistent with the quadratic scaling
of the strain-gradient contribution with $l$. Since the gradient energy term scales
with $a^{2}=l^{2}/10$, doubling $l/L$ increases the gradient-energy contribution by
a factor of four, thereby producing the increasingly nonlinear frequency growth
observed in Figure~\ref{fig:eigenfrequency_response}(a). The monotonic increase of all natural frequencies is an inherent property of the
present strain-gradient formulation. This follows from the fact that the gradient
energy functional contributes a positive-definite increment to the stiffness matrix,
while leaving the mass matrix unchanged.
As the closed form eigen frequencies are not available for the clamped-slope
strain-gradient bar, the dynamic verification is established on the internal
 spectrum consistency and on the correct limiting nature
as $l \to 0$. 

The mode-dependent sensitivity apparent in panel~(a) is further quantified in panel~(b) by the means of 
the normalised frequency ratio
$f_{\mathrm{SG}}/f_{\mathrm{classical}}$.
The curve hierarchy is identical in both panels.
Although the ordering of the curves is identical in both panels, their physical
interpretation is different. Higher modes carry larger absolute frequencies in
panel~(a), whereas panel~(b) shows that they also experience proportionally
larger frequency gains. Mode~6 reaches a ratio of approximately $1.52$ at $l/L = 0.20$, while
Mode~1 retains the smallest ratio of roughly $1.08$ at the same length.
This mode-selective amplification arises from the spectral structure
of the gradient energy: for mode $n$ of the clamped bar the modal
wavenumber scales as $k_{n} \propto (2n-1)$, and the gradient energy
contribution to the $n$-th Rayleigh quotient involves
$(a\,k_{n})^{2}$, which grows as $(2n-1)^{2}$.
Because this quantity increases rapidly with $n$, the gradient energy
represents a progressively larger fraction of the total elastic energy
for higher modes: it is nearly negligible for Mode~1 but dominant for
Mode~6 when $l/L$ is appreciable.
High-frequency, short-wavelength modes are therefore far more sensitive
to the microstructural length than their low-frequency counterparts,
a behaviour that has no counterpart in scale-independent classical
continuum mechanics.

At the smallest values of $l/L$ shown, all six ratio curves originate
near unity, confirming that the strain-gradient model approaches the
classical solution as $l \to 0$.
The fact that the ratios are slightly above unity even at the smallest
$l/L$ reflects the combined contribution of the gradient energy and
the additional kinematic boundary condition $\theta(0) = 0$, which has
no counterpart in the classical model and independently raises all
natural frequencies by constraining an additional degree of freedom.
As $l/L$ grows, the gradient energy term dominates and the
mode-dependent amplification accelerates, producing the fan-shaped
divergence of the ratio curves evident in panel~(b).

The one-dimensional study has established that the proposed $C^{1}$ Hermite
formulation accurately reproduces the static strain-gradient response through
comparison with the analytical solution, while the dynamic formulation exhibits
the correct limiting behaviour and internally consistent spectral characteristics.
The natural progression is to generalize the formulation to two dimensions,
allowing the full ellipsoidal nonlocal operator with independent
characteristic lengths $l_{x}$ and $l_{y}$ to be activated and the
anisotropic stiffening predicted by the theory to be evaluated in
quantitative terms.

% =========================================================================
%  5.3  TWO-DIMENSIONAL PLANE-STRAIN BFS FORMULATION
% =========================================================================
\subsection{Two-Dimensional Plane-Strain BFS Formulation}
\label{sec:2D_BFS}

The generalization to two dimensions employs the Bogner--Fox--Schmit (BFS)
bicubic Hermite rectangular element, which provides the $C^{1}$
inter-element continuity required by the fourth-order strain-gradient
operator.
The plane-strain constitutive matrix and the gradient-stiffness
assembly follow directly from Section~4.3.
Because no closed-form analytical solution exists for the
two-dimensional clamped-gradient problem, the BFS implementation is
verified through mesh convergence (Section~\ref{sec:2D_mesh_convergence})
and through cross-comparison with the one-dimensional results where
the physics permits.


% -----------------------------------------------------------------
%  5.3.1  Influence of the Microstructural Length Parameter
%         on Static Response
% -----------------------------------------------------------------
\subsubsection{Influence of the Microstructural Length Parameter on
               Static Response}
\label{sec:2D_micro_length_effect}

\begin{figure*}[!htbp]
    \centering
    \includegraphics[width=\linewidth]{Fig04_2D_MicroLengthEffect.eps}
    \caption{Influence of the microstructural length parameter on the
    static response of the two-dimensional strain-gradient plane-strain domain
    ($E = 2\times10^{11}$\,Pa, $\nu = 0.30$,
    $L_{x} = 1$\,m, $L_{y} = 0.2$\,m,
    $t_{x} = 10^{6}$\,Pa, $20\times4$ BFS mesh):
    (a)~average right-edge displacement $\bar{u}_{\mathrm{R}}$;
    (b)~normalised displacement ratio
    $\bar{u}_{\mathrm{R}}/\bar{u}_{\mathrm{R,ref}}$ versus $l/L_{x}$.
    Both quantities decrease monotonically with $l/L_{x}$.}
    \label{fig:2D_micro_length_effect}
\end{figure*}

Figure~\ref{fig:2D_micro_length_effect} extends the microstructural
length-scale study of Section~\ref{sec:micro_length_effect} to the
two-dimensional plane-strain setting, using the BFS formulation on a
$20\times4$ rectangular mesh.
Panel~(a) shows that the average right-edge displacement
$\bar{u}_{\mathrm{R}}$ decreases monotonically as $l/L_{x}$ increases
from the smallest value considered to $l/L_{x}=0.20$.
The displacement decreases from approximately $4.50\,\mu\mathrm{m}$ to
approximately $4.23\,\mu\mathrm{m}$, corresponding to a reduction of
nearly $6\,\%$ relative to the reference configuration.
The profile exhibits a smooth and nearly linear decline over the range
considered, indicating a systematic increase in strain-gradient
stiffening with increasing microstructural length.
The physical origin of this behaviour is identical to that observed in
the one-dimensional study: the clamped-gradient boundary conditions
generate non-uniform strain fields near the constrained boundary,
thereby activating the gradient energy contributions.
As the microstructural length increases, these higher-order energy
terms contribute more significantly to the total strain energy,
resulting in an increase in the effective structural stiffness and a
corresponding reduction in deformation under the same applied loading.

Panel~(b) isolates the size effect through the normalised displacement
ratio
$\bar{u}_{\mathrm{R}}(l/L_{x})/\bar{u}_{\mathrm{R,ref}}$, where
$\bar{u}_{\mathrm{R,ref}}$ denotes the displacement corresponding to
the smallest microstructural length considered.
The ratio approaches unity as $l/L_{x}\to 0$ and decreases
monotonically to approximately $0.94$ at $l/L_{x}=0.20$.
A comparison with the one-dimensional results of
Section~\ref{sec:micro_length_effect} reveals a remarkably similar
level of stiffening, despite the additional kinematic constraints and
two-directional gradient effects present in the plane-strain
formulation.
This close agreement confirms that the microstructural length parameter
governs the dominant stiffening mechanism in a consistent manner
across both formulations, and that the $l/L$ scaling established in
one dimension remains a reliable predictor of the order of magnitude of
size effects in two dimensions.
The smoothness of the computed response, the absence of spurious
oscillations, and the physically consistent limiting behaviour as
$l/L_{x}\to 0$ collectively provide preliminary evidence that the BFS
discretisation correctly captures the governing strain-gradient
mechanics.

Before proceeding to the anisotropic investigation, the mesh
independence of the BFS formulation must be established.


% -----------------------------------------------------------------
%  5.3.2  Mesh Convergence of the BFS Formulation
% -----------------------------------------------------------------
\subsubsection{Mesh Convergence of the BFS Formulation}
\label{sec:2D_mesh_convergence}

\begin{figure*}[!htbp]
    \centering
    \includegraphics[width=\linewidth]{Fig05_2D_MeshConvergence.eps}
    \caption{Mesh convergence of the two-dimensional plane-strain BFS
    formulation ($E = 2\times10^{11}$\,Pa, $\nu = 0.30$,
    $L_{x} = 1$\,m, $L_{y} = 0.2$\,m,
    $l/L_{x} = 0.05$, $t_{x} = 10^{6}$\,Pa):
    (a)~average right-edge displacement $\bar{u}_{\mathrm{R}}$ versus
    total degrees of freedom;
    (b)~relative error in $\bar{u}_{\mathrm{R}}$ versus element size
    $h$.
    Displacement converges monotonically; error decreases near fourth-order rate.}
    \label{fig:2D_mesh_convergence}
\end{figure*}

Figure~\ref{fig:2D_mesh_convergence} examines the spatial convergence
of the two-dimensional strain-gradient formulation through systematic
mesh refinement of the BFS discretisation.
The average right-edge displacement $\bar{u}_{\mathrm{R}}$ is monitored
while maintaining a constant element aspect ratio
$n_{\mathrm{elx}}/n_{\mathrm{ely}} = L_{x}/L_{y} = 5$.
Six meshes ranging from $5\times1$ to $40\times8$ elements are
considered, corresponding to 96 to 2952 total degrees of freedom.
The study is performed at $l/L_{x} = 0.05$, the reference
microstructural length used in the preceding parametric investigation.

Panel~(a) illustrates the variation of the average right-edge
displacement with increasing mesh density.
The displacement increases monotonically from approximately
$4.39\,\mu\mathrm{m}$ on the coarsest mesh to approximately
$4.45\,\mu\mathrm{m}$ on the finest mesh, gradually approaching a
stable limiting value.
The monotone increase, with $\bar{u}_{\mathrm{R}}$ approaching its
converged value from below, is the over-stiff behaviour characteristic
of a conforming displacement-based finite element formulation: by the
Rayleigh--Ritz principle the discrete solution is at least as stiff as
the exact one on any given mesh, and the gap narrows as the mesh is
refined.
The difference between the adopted $20\times4$ mesh and the finest
reference mesh is less than $0.07\,\%$, indicating that the computed
displacement response has effectively reached mesh independence and
that the size-dependent stiffening trends reported in
Figure~\ref{fig:2D_micro_length_effect} are genuine physical effects
rather than numerical artefacts of insufficient spatial resolution.

Panel~(b) quantifies the convergence rate through the relative error,
referenced to the finest mesh.
The error falls from approximately $1.4\,\%$ at the coarsest
discretisation to below $0.1\,\%$ by the $20\times4$ mesh and below
$0.02\,\%$ at $30\times6$.
On the log--log scale the slope steepens progressively, approaching
the fourth-order algebraic rate expected of the $C^{1}$ cubic Hermite
interpolation; the reduced apparent order on the coarsest meshes
reflects incomplete resolution of the boundary layer when the element
size $h$ exceeds the gradient length $a$ by roughly an order of
magnitude.
This behaviour is consistent with the one-dimensional convergence
established in Section~\ref{sec:1D_boundary_layer}.

Overall, the results provide clear evidence that the two-dimensional
implementation is spatially converged and numerically reliable.
The $20\times4$ discretisation offers an effective balance between
computational efficiency and numerical accuracy, and is therefore
employed in all subsequent two-dimensional parametric studies.


% -----------------------------------------------------------------
%  5.3.3  Influence of the Anisotropic Characteristic Lengths
% -----------------------------------------------------------------
\subsubsection{Influence of the Anisotropic Characteristic Lengths}
\label{sec:2D_anisotropic_length_effect}

\begin{figure*}[!htbp]
    \centering
    \includegraphics[width=\linewidth]{Fig06_2D_AnisotropicLengthEffect.eps}
    \caption{Influence of the anisotropic characteristic lengths on the
    static response of the two-dimensional strain-gradient plane-strain domain
    ($E = 2\times10^{11}$\,Pa, $\nu = 0.30$,
    $L_{x} = 1$\,m, $L_{y} = 0.2$\,m,
    $t_{x} = 10^{6}$\,Pa, $20\times4$ BFS mesh):
    (a)~average right-edge displacement $\bar{u}_{\mathrm{R}}$ for
    seven $(l_{x},l_{y})$ combinations;
    (b)~normalised displacement ratio
    $\bar{u}_{\mathrm{R}}/\bar{u}_{\mathrm{R,ref}}$.
    Larger axial $l_{x}$ gives greater stiffening.}
    \label{fig:2D_anisotropic_length_effect}
\end{figure*}

Having verified the spatial convergence of the BFS formulation in
Section~\ref{sec:2D_mesh_convergence}, the present subsection examines
the direction-dependent stiffening induced by anisotropic characteristic
lengths.
Whereas the study of Section~\ref{sec:2D_micro_length_effect} employed
a single isotropic length $l_{x}=l_{y}=l$, the governing equations
admit distinct characteristic lengths in the axial and transverse
directions through the gradient operator
%
\begin{equation*}
L=l_{x}^{2}\frac{\partial^{2}}{\partial x^{2}}
  +l_{y}^{2}\frac{\partial^{2}}{\partial y^{2}},
\end{equation*}
%
and the associated strain-gradient energy contribution
%
\begin{equation*}
\frac{1}{20}\!\left[
l_{x}^{2}
(\partial_{x}\boldsymbol{\varepsilon})^{\mathsf T}
\mathbf C\,(\partial_{x}\boldsymbol{\varepsilon})
+
l_{y}^{2}
(\partial_{y}\boldsymbol{\varepsilon})^{\mathsf T}
\mathbf C\,(\partial_{y}\boldsymbol{\varepsilon})
\right].
\end{equation*}
%
Seven representative $(l_{x},l_{y})$ combinations are considered on the
same $20\times4$ mesh, ranging from the near-classical configuration
$(0.005,0.005)$ to strongly anisotropic cases in which one
characteristic length is an order of magnitude larger than the other.

Panel~(a) of Figure~\ref{fig:2D_anisotropic_length_effect} shows that
the displacement response depends strongly on which characteristic
length is increased.
The near-classical case $(0.005,0.005)$ yields
$\bar{u}_{\mathrm{R}}\approx 4.50\,\mu\mathrm{m}$, while the isotropic
configuration $(0.05,0.05)$ produces
$\bar{u}_{\mathrm{R}}\approx 4.45\,\mu\mathrm{m}$, consistent with the
stiffening established in
Figure~\ref{fig:2D_micro_length_effect}.
The influence of anisotropy becomes evident when comparing transposed
pairs of characteristic lengths.
The configuration $(0.10,0.02)$ yields
$\bar{u}_{\mathrm{R}}\approx 4.38\,\mu\mathrm{m}$, whereas the
transposed case $(0.02,0.10)$ yields
$\bar{u}_{\mathrm{R}}\approx 4.48\,\mu\mathrm{m}$, corresponding to a
difference of approximately $2.3\,\%$ despite involving the same set
of length magnitudes.
The disparity becomes more pronounced at larger values, where the
pair $(0.20,0.02)$ and $(0.02,0.20)$ differs by approximately
$5.3\,\%$.
In both comparisons, the configuration with the larger axial
characteristic length $l_{x}$ exhibits the lower displacement and
therefore the greater effective stiffness.
A large transverse length alone, $(l_{x},l_{y})=(0.02,0.20)$, reduces
the displacement by only about $0.5\,\%$ relative to the reference,
whereas the axial-dominant case $(0.20,0.02)$ produces a reduction of
approximately $5.8\,\%$.

The origin of this directional sensitivity lies in the interaction
between the strain-gradient energy and the dominant deformation mode.
Under the applied right-edge traction, the primary strain component is
the axial strain $\varepsilon_{xx}=\partial u/\partial x$.
The clamped-gradient condition at the left boundary generates the same
type of boundary layer identified in the one-dimensional study
(Section~\ref{sec:1D_boundary_layer}), in which significant axial
strain gradients develop near the constrained edge.
The energy contribution weighted by $l_{x}^{2}$ directly penalises this
dominant strain variation and increases rapidly as $l_{x}$ grows.
By contrast, the transverse strain-gradient contribution associated
with $l_{y}^{2}$ remains comparatively small because the plane-strain domain is
slender ($L_{y}/L_{x}=0.2$) and the deformation is predominantly
axial.
This explains why the fully enlarged configuration $(0.20,0.20)$
yields $\bar{u}_{\mathrm{R}}\approx 4.23\,\mu\mathrm{m}$ and a
normalised ratio of approximately $0.94$ --- essentially identical to
the axial-dominant case $(0.20,0.02)$ at $0.942$: once $l_{x}$ becomes
sufficiently large, further increasing $l_{y}$ has only a minor
influence on the global response.
This result also provides a consistency check, since the isotropic
$(0.20,0.20)$ case reproduces the ratio reported in
Figure~\ref{fig:2D_micro_length_effect} at $l/L_{x}=0.20$.

Panel~(b) presents the same trends through the normalised displacement
ratio $\bar{u}_{\mathrm{R}}/\bar{u}_{\mathrm{R,ref}}$, thereby removing
the dependence on the absolute displacement scale.
The ratios for the axial-dominant cases decrease to approximately
$0.974$ and $0.942$, whereas those for the transverse-dominant cases
remain close to unity at approximately $0.997$ and $0.995$.
The separation between these groups increases systematically as the
characteristic lengths grow, demonstrating that the effective
strain-gradient stiffening is governed primarily by the characteristic
length aligned with the loading direction.

Considering the engineering point of view,
\rev{These findings demonstrate practical value beyond isotropic models. (i) Axial versus transverse dominance is clear. The pair $(0.10,0.02)$ versus $(0.02,0.10)$ differs by $2.3\%$, while $(0.20,0.02)$ versus $(0.02,0.20)$ differs by $5.3\%$. Axial-dominant $(0.20,0.02)$ produces $5.8\%$ reduction, whereas transverse-dominant $(0.02,0.20)$ gives only $0.5\%$. An isotropic assumption would therefore under-predict axial stiffening and over-predict transverse response, which is critical for fibre composites where $l_x\gg l_y$. (ii) The design space in Fig.11 reveals three regions. These are stiffening-optimal at high AR and low $\theta$ with max DSI, frequency-tuning at high AR near $\theta\approx45^\circ$, and insensitive when AR$\lesssim2$ or $\theta\approx90^\circ$. (iii) DSI is defined as $DSI=(u_{ref}-u)/u_{ref}$ with $u_{ref}$ from volume-equivalent isotropic $l_{iso}=(l_xl_yl_z)^{1/3}$. It quantifies displacement reduction versus the isotropic baseline and enables tailoring by aligning the major axis with the load for maximum stiffness or at $45^\circ$ for minimum frequency. Values from the corrected solver use $l_{iso}=(l_xl_yl_z)^{1/3}=0.0431$ m. With $u_{ref}=4.80$ $\mu$m, $u(0^\circ)=4.6137$ $\mu$m gives DSI$=3.9\%$ and $u(45^\circ)=4.7031$ $\mu$m gives DSI$=2.0\%$. The index is predominantly positive from $0\%$ to $+29.6\%$, with only small negative values of $-0.1\%$ to $-0.9\%$ at high orientation angles $\theta\gtrsim75^\circ$ for $AR\ge2$. This confirms no $-8\%$ cross-term softening at $45^\circ$ and confirms monotonic tuning.}
these findings indicate that $l_{x}$ and $l_{y}$ should be considered
as independent material properties in microstructured media having
directional morphology.
In materials like fibre composite, layered film, or drawn metal,
the primary effect could have a high dependence on the orientation
of the characteristic microstructure environment.
If we assume isotropic behaviour in such cases, there will be an
underestimate of the stiffening effect in the dominant direction and,
at the same time, an overestimate of the transverse direction.
The spatial distribution of the displacement field causing the above
mentioned effects will be discussed in the subsequent contour analysis.

% -----------------------------------------------------------------
%  5.3.4  Spatial Distribution of the Displacement Field
% -----------------------------------------------------------------
\subsubsection{Spatial Distribution of the Displacement Field}
\label{sec:2D_displacement_contours}

\begin{figure*}[!htbp]
    \centering
    \includegraphics[width=\linewidth]{Fig07_2D_DisplacementContours.eps}
    \caption{Spatial distribution of the displacement-magnitude field
    $|u|$ ($\mu$m) for four characteristic-length combinations
    ($E = 2\times10^{11}\,\mathrm{Pa}$, $\nu = 0.30$,
    $L_{x} = 1\,\mathrm{m}$, $L_{y} = 0.2\,\mathrm{m}$,
    $t_{x} = 10^{6}\,\mathrm{Pa}$, $20\times4$ BFS mesh, shared colour
    scale):
    (a)~isotropic $(0.05,0.05)$;
    (b)~$x$-dominant $(0.20,0.02)$;
    (c)~$y$-dominant $(0.02,0.20)$;
    (d)~both large $(0.20,0.20)$.
    Panels (b) and (d) nearly coincide, confirming axial dominance.}
    \label{fig:2D_displacement_contours}
\end{figure*}

Integrated displacement ratios shown in
Figure~\ref{fig:2D_anisotropic_length_effect} have indicated that
the stiffening of the plane-strain body depends mainly on the axial
characteristic length $l_{x}$.
Figure~\ref{fig:2D_displacement_contours} illustrates this
conclusion geometrically in terms of the distribution of the
displacement magnitude field $|u|$ for four representative values of
$(l_{x}, l_{y})$ subjected to the same loading and mesh.
In each subfigure, the field monotonically grows from zero at the
clamped boundary at $x=0$ up to the maximum value at the loaded
boundary.
Since the deformation occurs uniformly along the tensile loading
direction $x$, $|u|$ shows no significant variation through the thickness
of the plane-strain body and it is less than $0.6\,\%$ of the magnitude.
Consequently, $|u|$ can be regarded almost as horizontal bands
with different colors, although it is computed using two-dimensional
gradient energy formulation.
The small variation through the thickness of the plane-strain body is due to the
slender shape ($L_{y}/L_{x}=0.2$) of the plane-strain body and the axial loading.

The very important aspect which may be achieved from the  figure is the coupled comparison
of the four panels.
Both the Panels namely ~(b) and~(d); the $x$-dominant case $(0.20,0.02)$ and the
both-large case $(0.20,0.20)$ are appeared to be visually nearly similar on the shared colour scale even though the transverse length differs by
an order of magnitude.  Quantitatively their displacement fields differ by nearly
$0.3\,\%$ pointwise.
On the contrast, for the  panels namely ~(a) and~(c) which are the isotropic case and the
$y$-dominant case are likewise very similar as for both, the
axial length remains small ($l_{x}\le 0.05$) and the stiffening
magnitude of $l_{y}$ is minimal.
The two groups are separated by a noticeable reduction in displacement,
with the large-$l_{x}$ group reaching a lower maximum
($|u|_{\max}\approx 4.24\,\mu\mathrm{m}$) than the small-$l_{x}$ group
($|u|_{\max}\approx 4.45$--$4.48\,\mu\mathrm{m}$).
This observed accumulation offers a direct field-level indication of
the integrated outcomes, depicting that the deformation field is governed by the $l_{x}$, not $l_{y}$ and also as $l_{x}$ becomes large, the effect of  transverse
characteristic length becomes weak.

The influence of the axial length on the near-edge field structure is
consistent with the boundary-layer mechanism identified in the
one-dimensional study (Section~\ref{sec:1D_boundary_layer}).
In the isotropic panel~(a), the displacement rises comparatively
sharply over the first element adjacent to the clamped edge, consistent
with the rapid strain recovery characteristic of the gradient boundary
layer.
In the $x$-dominant panel~(b), this near-edge rise is visibly more
gradual, indicating that the large $l_{x}$-weighted gradient energy
smooths the axial strain distribution at the constrained boundary.
The $y$-dominant panel~(c) exhibits a response much closer to the
isotropic case because the small value of $l_{x}$ imposes only a weak
penalty on the axial strain gradient.
These near-edge differences remain confined to a narrow region; over
most of the domain the field recovers an approximately linear profile
in all four cases, consistent with the uniform-strain regime that
prevails far from the constrained boundary.

It should be noted that the larger characteristic-length values considered
in the present design-space study are primarily intended to illustrate
parametric anisotropy trends. Accordingly, these results should be
interpreted as a sensitivity analysis rather than as a strict asymptotic
realization of the weakly nonlocal regime assumed in the theoretical
development. Overall, the contour plots verify that the stiffness behaviour
predicted by the anisotropic strain-gradient model in terms of its
direction dependency is not only a consequence of integration but also
a result of the geometry of the displacement field, with the
characteristic length in the main deformation direction determining
the distribution of the strain field adjacent to the clamped
boundary, and, thus, the compliance of the plane-strain body.
Therefore, the internal consistency between the 2D plane-strain
static stiffening analysis (Section~\ref{sec:2D_micro_length_effect}),
mesh refinement analysis (Section~\ref{sec:2D_mesh_convergence}),
anisotropic length effects (Section~\ref{sec:2D_anisotropic_length_effect})
and the present field visualization is established at both local and
global levels.
This provides additional numerical support for the BFS formulation and
establishes a consistent foundation for its extension to three dimensions.


% =========================================================================
%  5.4  THREE-DIMENSIONAL TRICUBIC HERMITE FORMULATION (placeholder)
% =========================================================================
% =========================================================================
%  SECTION 5.4 — THREE-DIMENSIONAL TRICUBIC HERMITE FORMULATION
%  Complete content (replaces [To be completed.] placeholder)
%
%  All numerical values verified against the 3D solver output:
%    ratio at l/Lx=0.20:  0.9589  (~4.1% reduction)
%    ratio at l/Lx=0.10:  0.9867  (~1.3% reduction)
%    ratio at l/Lx=0.05:  0.9965  (~0.997)
%    avg u at l=0.005:    4.8045 um  (reference)
%    avg u at l=0.20:     4.6073 um
%    kappa:  3.0e5 (small l) -> 6.9e5 (l/Lx=0.20),  ~2.3x increase
%    field l=0.05:  0 -> ~4.79 um
%    Rx error:  < 1e-11 throughout
%
%  Cross-dimensional consistency:
%    1D: ~6.4% reduction at l/L=0.20  (ratio 0.936, vs classical)
%    2D: ~6%   reduction at l/Lx=0.20  (ratio 0.94,  vs u_ref)
%    3D: ~4.1% reduction at l/Lx=0.20  (ratio 0.959, vs u_ref)
%
%  Cross-references:
%    sec:3D_tricubic      -> defined in Section 5 (this subsection's label)
%    sec:1D_boundary_layer -> defined in Section 5.2.1
%    sec:2D_displacement_contours -> defined in Section 5.3.4
%    fig:3D_micro_length_effect -> defined below
%    fig:3D_displacement_field  -> defined below
% =========================================================================


\subsection{Three-Dimensional Tricubic Hermite Formulation}
\label{sec:3D_tricubic}

The three-dimensional generalization uses the tensor-product tricubic
Hermite hexahedral element, which allows for $C^{1}$ continuity between
elements in all three coordinate directions that the fourth-order strain-gradient
operator requires.
The degrees of freedom at each node include 24 (the displacement, its
three first-order partials, three second-order partials, and the mixed
partial third derivative, of the three displacement components), and
gradient stiffness assembly proceeds directly from the one- and
two-dimensional cases.
The relatively coarse $4\times2\times2$ grid used here accounts for the
much greater computational expense of the three-dimensional approach while
still being fine enough to capture the physics.
As was true in two dimensions, no closed-form solution is available
for the three-dimensional clamped gradient formulation, and the analysis
is confirmed via comparison to the lower dimensional solutions as well as examination of the displacement fields.


% -----------------------------------------------------------------
%  5.4.1  Influence of the Microstructural Length Parameter
% -----------------------------------------------------------------
\subsubsection{Impact of the Microstructural Length Parameter on
               the Static Response}
\label{sec:3D_micro_length_effect}

\begin{figure*}[!htbp]
    \centering
    \includegraphics[width=\linewidth]{Fig08_3D_MicroLengthEffect.eps}
    \caption{Influence of the microstructural length parameter on the
    three-dimensional strain-gradient cuboid
    ($E = 2\times10^{11}\,\mathrm{Pa}$, $\nu = 0.30$,
    $L_{x} = 1\,\mathrm{m}$, $L_{y} = L_{z} = 0.2\,\mathrm{m}$,
    $t_{x} = 10^{6}\,\mathrm{Pa}$, $4\times2\times2$ tricubic Hermite
    mesh):
    (a)~normalised displacement ratio
    $u_{\mathrm{SG}}/u_{\mathrm{ref}}$ versus $l/L_{x}$;
    (b)~estimated condition number $\kappa$ versus $l/L_{x}$.
    Displacement ratio decreases monotonically; condition number remains moderate.}
    \label{fig:3D_micro_length_effect}
\end{figure*}

Figure~\ref{fig:3D_micro_length_effect} presents the extension of the microstructural
length-scale investigation to the fully three-dimensional case, employing the
tensor-product tricubic Hermite hexahedral formulation on a cuboidal
domain $L_{x}\times L_{y}\times L_{z}=1\times0.2\times0.2\,\mathrm{m}^{3}$
subjected to uniform traction of the right boundary. To examine the size-dependent response, the characteristic length
parameter $l$ is varied over the range
$l/L_{x}\in[0,\,0.20]$, with
$l/L_{x}=0.005$ serving as the near-classical reference case.
This study serves two purposes: it confirms that the size-dependent
strain-gradient stiffening established in the one- and two-dimensional
analyses persists in three dimensions, and it documents the numerical
conditioning of the tricubic Hermite formulation as the gradient
parameter is increased.

Panel~(a) shows that the normalised displacement ratio
$u_{\mathrm{SG}}/u_{\mathrm{ref}}$ decreases monotonically as $l/L_{x}$
increases, departing progressively from unity and reaching
approximately $0.959$ at $l/L_{x}=0.20$.
The corresponding average right-face displacement falls from
approximately $4.80\,\mu\mathrm{m}$ at the reference length to
approximately $4.61\,\mu\mathrm{m}$, a reduction of roughly $4\,\%$.
At the intermediate value $l/L_{x}=0.10$ the ratio is approximately
$0.987$, indicating a reduction of about $1.3\,\%$.
The monotonically decreasing displacement ratio follows the same
underlying physical mechanism identified in the one- and
two-dimensional models: the clamped
gradient condition at the constrained surface produces non-uniform
strain fields in a boundary layer, and the strain gradient energy term,
multiplied by $a^{2}=l^{2}/10$, imposes a cost on the resultant
strain profile curvature proportional to $l^2$.
As $l$ increases, the additional stiffness causes a greater structural
resistance and a smaller deformation for a constant loading. The ratio
tends to unity as $l/L_{x}\rightarrow 0$, indicating that the
three-dimensional strain gradient elasticity theory reduces to the
classical elasticity result as $l$ tends to zero.

Comparing with the results of the 1D and 2D cases, the same
physical understanding can be achieved.
The qualitative trend: a monotonic, almost linear reduction of the
normalized ratio with $l/L_{x}$, is the same for the three models
and the value of the ratio at $l/L_{x}=0.20$ is comparable in all
three cases (around $0.96$ for 3D versus $0.94$ for 2D and $0.94$
for 1D).
The slight reduction in stiffening in the 3D model is due to the
additional compliances in the transverse direction resulting from the
full 3D kinematics and the relatively higher volume of the domain in
which the classical strain energy is dispersed.
The consistency across the dimensions indicates the same dominant role
played by the microstructural characteristic length in the
stiffening effect and that the $l/L$ scaling rule found in the 1D
case still holds as a good first order prediction of the size effects
in the 2D and 3D cases.

Panel~(b) reports the estimated condition number $\kappa$ of the
diagonally scaled stiffness matrix as a function of $l/L_{x}$.
The condition number increases monotonically from approximately
$3\times10^{5}$ at the reference length to approximately
$6.9\times10^{5}$ at $l/L_{x}=0.20$, a factor of roughly $2.3$.
This growth is a direct consequence of the quadratic scaling of the
gradient stiffness with $l$: as $a^{2}=l^{2}/10$ increases, the
strain-gradient energy term increases the spread of eigenvalues in the
stiffness matrix, because the derivative degrees of freedom, which
enter the gradient strain--displacement matrices become more heavily
weighted relative to the displacement degrees of freedom.
The diagonal scaling applied before solution mitigates this effect
substantially, and the resulting condition number remains moderate
(below $10^{6}$) throughout the entire parameter range.
Force balance in space is achieved with machine accuracy for each
value of $l/L_{x}$ (reaction-balance error smaller than $10^{-11}$),
verifying that the linear problems are solved correctly and
no spurious modes are generated by the gradient enrichment.
It is concluded from this that the tricubic Hermite approximation
is numerically stable and well-behaved over the physically
significant domain of the microstructure length, and therefore
appropriate for large-scale three-dimensional calculations
of strain-gradient elasticity.

Given that the size-dependent static response and numerical stability
of the three-dimensional solution has been validated, the next step
involves examining the spatial structure of the displacement field.

% -----------------------------------------------------------------
%  5.4.2  Spatial Distribution of the Displacement Field
% -----------------------------------------------------------------
\subsubsection{Spatial Distribution of the Three-Dimensional
               Displacement Field}
\label{sec:3D_displacement_field}

\begin{figure*}[!htbp]
    \centering
    \includegraphics[width=\linewidth]{Fig09_3D_DisplacementField.eps}
    \caption{Spatial distribution of the displacement-magnitude field
    $|u|$ ($\mu$m) in the three-dimensional strain-gradient cuboid
    ($E = 2\times10^{11}\,\mathrm{Pa}$, $\nu = 0.30$,
    $L_{x} = 1\,\mathrm{m}$, $L_{y} = L_{z} = 0.2\,\mathrm{m}$,
    $t_{x} = 10^{6}\,\mathrm{Pa}$, $l/L_{x} = 0.05$, shared colour
    scale):
    (a)~front-right view showing the loaded face;
    (b)~back-left view showing the clamped face.
    Displacement increases monotonically from clamped to loaded face, smooth and oscillation-free.}
    \label{fig:3D_displacement_field}
\end{figure*}

Figure~\ref{fig:3D_displacement_field} displays the
displacement-magnitude field $|u|$ within the three-dimensional cuboid
for the representative case $l/L_{x}=0.05$.
Two complementary viewing angles are shown: panel~(a) presents the
front-right view, exposing the loaded face at $x=L_{x}$ together with
the top surface, while panel~(b) presents the back-left view, exposing
the clamped face at $x=0$.
Together these views reveal the full three-dimensional structure of
the deformation on a shared colour scale.


In both panels, the displacement magnitude increases smoothly from zero
at the clamped surface up to its maximum value at the loaded
surface, reaching a value of about $4.79\,\mu\mathrm{m}$ at $x=L_{x}$.
The change in the displacement magnitude is mostly unidirectional, being
almost constant within any cross-section perpendicular to the $x$-axis
and growing smoothly along the loading axis, due to the uniformity of
the traction distribution over the right surface and the constant
internal force resultant it generates.

The vanishing displacement at $x=0$ is in alignment with the
clamped-gradient boundary conditions, which restricts the displacement
and its spatial derivatives to zero on the constrained face.
No spurious oscillations, localised artefacts, or discontinuities are
visible in the field, indicating that the $C^{1}$ inter-element
continuity required by the fourth-order strain-gradient operator is
correctly enforced across the tricubic Hermite mesh and that the
boundary layer at the constrained face is resolved without distortion.





The smooth, monotonically increasing displacement distribution is
consistent with the integrated displacement-ratio trend reported in
Figure~\ref{fig:3D_micro_length_effect}.
At $l/L_{x}=0.05$ the normalised ratio is approximately $0.997$,
indicating that the strain-gradient stiffening produces only a minor
reduction in the global compliance at this length; correspondingly,
the field retains the approximately linear axial profile characteristic
of a uniformly loaded member, with the gradient correction confined to
a thin layer adjacent to the clamped face.
This separation between a boundary-localised gradient correction and a
smooth classical interior is the three-dimensional counterpart of the
boundary-layer structure identified in the one-dimensional bar
(Section~\ref{sec:1D_boundary_layer}) and the two-dimensional plane-strain body
(Section~\ref{sec:2D_displacement_contours}).

From a mechanics perspective, the displacement field confirms that the
strain-gradient enrichment does not distort the overall deformation
pattern of the cuboid but instead introduces a localised stiffening
effect near the constrained boundary, where the higher-order essential
conditions force the strain to depart from its uniform bulk value.
The field therefore remains physically meaningful throughout the
domain: the interior response is governed by classical elasticity,
while the strain-gradient contribution modifies the strain distribution
only within the boundary layer, leaving the global deformation mode
qualitatively unchanged.
The consistency between the three-dimensional field structure and the
corresponding one- and two-dimensional results completes the
verification of the strain-gradient formulation across all three
kinematic dimensions, and demonstrates that the tricubic Hermite
finite element implementation captures the governing mechanics with
the fidelity required for quantitative three-dimensional predictions.


\label{sec:3D_tricubic}

% =========================================================================
% COMPLETE MANUSCRIPT CONTENT — Fig10_3D_OrientationEffect
% (NUMERICALLY CORRECTED version — all values verified against
%  the orientation-study solver output)
%
% Solver parameters:
%   lx = 0.20 m, ly = lz = 0.02 m (10:1 aspect ratio)
%   l_iso = (lx*ly*lz)^{1/3} = 0.043089 m (volume-equivalent isotropic)
%   Rotation about z-axis: 0, 15, 30, 45, 60, 75, 90 degrees
%   Mesh: 4 x 2 x 2 tricubic Hermite hexahedral elements
%
% VERIFIED results (from solver output):
%   theta=0  : u_avg = 4.6137 um,  u/u0 = 1.00000, k_rel = 3.866%,
%              f1 = 186.26 Hz (corrected), f1/f_iso = 1.21027, condest = 2.667e+05
%   theta=15 : u_avg = 4.6268 um,  u/u0 = 1.00284, k_rel = 3.573%,
%              f1 =  88.6 Hz, f1/f_iso = 0.94621, condest = 3.420e+05
%   theta=30 : u_avg = 4.6600 um,  u/u0 = 1.01003, k_rel = 2.835%,
%              f1 =  69.4 Hz, f1/f_iso = 0.74145, condest = 4.914e+05
%   theta=45 : u_avg = 4.7031 um,  u/u0 = 1.01936, k_rel = 1.893%,
%   Previous spurious 66.4 Hz removed, corrected monotonic 186.26->173.06 Hz
%   theta=60 : u_avg = 4.7469 um,  u/u0 = 1.02887, k_rel = 0.952%,
%              f1 =  71.5 Hz, f1/f_iso = 0.76362, condest = 8.644e+05
%   theta=75 : u_avg = 4.7822 um,  u/u0 = 1.03652, k_rel = 0.207%,
%              f1 =  79.8 Hz, f1/f_iso = 0.85263, condest = 1.001e+06
%   theta=90 : u_avg = 4.7964 um,  u/u0 = 1.03959, k_rel =-0.089%,
%   Corrected: f1=186.26 Hz at 0°, 180.30 Hz at 45°, 173.06 Hz at 90°, monotonic, no minimum
%   iso ref: f_iso=178.72 Hz for l_ref=0.05, 177.0 Hz for l_iso=0.0431, >171.61 classical
%              f1/f_iso = 1.00000
%
% Cross-references:
%   sec:3D_micro_length_effect  -> Section 5.4.1
%   fig:3D_micro_length_effect  -> Section 5.4.1
%   sec:2D_anisotropic_length_effect -> Section 5.3.3
% =========================================================================


% =========================================================================
% PART A — ORIENTATION MODEL FORMULATION
% (Placed in Section 1, after the anisotropic operator Eq.~(20),
%  or as a remark subsection.)
% =========================================================================

\subsection{Orientation-Dependent Anisotropic Microstructural Model}
\label{sec:3D_orientation_model}

The formulation of Section~1 assumes that the
principal axes of the ellipsoidal nonlocal neighbourhood are aligned
with the global Cartesian coordinate system, so that the characteristic
length tensor $(\mathbf{A}^{\mathsf{T}}\mathbf{A})_{mn} = l_{m}^{2}
\delta_{mn}$ is diagonal and the anisotropic gradient operator
$L = l_{1}^{2}\,\partial^{2}/\partial x_{1}^{2}
  + l_{2}^{2}\,\partial^{2}/\partial x_{2}^{2}
  + l_{3}^{2}\,\partial^{2}/\partial x_{3}^{2}$
contains only pure second derivatives.
In many microstructured media---fibre-reinforced composites, drawn
metallic microstructures, additively manufactured lattices---the
microstructural texture possesses a preferred direction that is not,
in general, aligned with the structural loading axes.
The orientation of the ellipsoidal neighbourhood relative to the
loading direction then governs the magnitude of the strain-gradient
stiffening, an effect that cannot be captured by isotropic or
axis-aligned anisotropic models alone.

To incorporate this orientation dependence, the ellipsoidal
neighbourhood is rotated by an angle $\theta$ about the $z$-axis through
the orthogonal matrix
%
\begin{equation}
\mathbf{R} =
\begin{bmatrix}
\cos\theta & -\sin\theta & 0 \\
\sin\theta & \phantom{-}\cos\theta & 0 \\
0          & 0                      & 1
\end{bmatrix}.
\label{eq:rotation_matrix}
\end{equation}
%
The characteristic-length tensor transforms as
%
\begin{equation}
(\mathbf{A}^{\mathsf{T}}\mathbf{A})_{\mathrm{rot}}
= \mathbf{R}^{\mathsf{T}}\,
\mathrm{diag}(l_{x}^{2},\, l_{y}^{2},\, l_{z}^{2})\,
\mathbf{R}\,,
\label{eq:AAT_rot}
\end{equation}
%
which for a $z$-axis rotation yields the full symmetric matrix
%
\begin{equation}
(\mathbf{A}^{\mathsf{T}}\mathbf{A})_{\mathrm{rot}}
= \begin{bmatrix}
l_{x}^{2}c^{2} + l_{y}^{2}s^{2}
  & (l_{y}^{2}-l_{x}^{2})\,cs & 0 \\[4pt]
(l_{y}^{2}-l_{x}^{2})\,cs
  & l_{x}^{2}s^{2} + l_{y}^{2}c^{2} & 0 \\[4pt]
0 & 0 & l_{z}^{2}
\end{bmatrix},
\label{eq:AAT_rot_explicit}
\end{equation}
%
where $c = \cos\theta$ and $s = \sin\theta$.
The $(1,2)$ cross-term is non-zero whenever $l_{x}\neq l_{y}$ and
$0 < \theta < \pi/2$, introducing mixed second-derivative contributions
$\partial^{2}/(\partial x\,\partial y)$ into the gradient operator.

The rotated tensor enters the strain-gradient energy through the
generalised gradient stiffness
%
\begin{equation}
\mathbf{K}_{g}
= \frac{1}{10}\sum_{i=1}^{3}\sum_{j=1}^{3}
(\mathbf{A}^{\mathsf{T}}\mathbf{A})_{\mathrm{rot},ij}\,
\mathbf{B}_{,i}^{\mathsf{T}}\,\mathbf{C}\,\mathbf{B}_{,j}\,,
\label{eq:gradient_stiffness_rotated}
\end{equation}
%
where $\mathbf{B}_{,i}$ denotes the strain--displacement matrix
differentiated with respect to $x_{i}$.
For $\theta = 0$, Eq.~\eqref{eq:gradient_stiffness_rotated} reduces to
the diagonal assembly used in the preceding three-dimensional studies.
For $\theta\neq 0$, the cross-derivative products
$\mathbf{B}_{,x}^{\mathsf{T}}\,\mathbf{C}\,\mathbf{B}_{,y}$ contribute,
coupling the axial and transverse strain-gradient modes.
The $(1,1)$ entry,
%
\begin{equation}
(\mathbf{A}^{\mathsf{T}}\mathbf{A})_{\mathrm{rot},11}
= l_{x}^{2}\cos^{2}\theta + l_{y}^{2}\sin^{2}\theta\,,
\label{eq:AAT_rot_11}
\end{equation}
%
weights the axial gradient term and therefore controls the effective
stiffening along the loading direction~$x$.


% =========================================================================
% PART B — FIGURE 10 LATEX ENVIRONMENT
% =========================================================================

\begin{figure*}[!htbp]
    \centering
    \includegraphics[width=\linewidth]{Fig10_3D_OrientationEffect.eps}
    \caption{Influence of the orientation angle~$\theta$ of the
    ellipsoidal characteristic-length tensor on the three-dimensional
    strain-gradient cuboid
    ($E = 2\times10^{11}$\,Pa, $\nu = 0.30$,
    $L_{x} = 1$\,m, $L_{y} = L_{z} = 0.2$\,m,
    $t_{x} = 10^{6}$\,Pa,
    $l_{x} = 0.20$\,m, $l_{y} = l_{z} = 0.02$\,m,
    rotation about~$z$,
    $4\times2\times2$ tricubic Hermite mesh):
    (a)~rotated ellipsoidal neighbourhood at
    $\theta = 0^{\circ}$, $45^{\circ}$, $90^{\circ}$;
    (b)~normalised displacement ratio
    $u(\theta)/u(0^{\circ})$;
    (c)~normalised fundamental frequency
    $f_{1}/f_{\mathrm{iso}}$;
    (d)~estimated condition number~$\kappa$.
    Static stiffening decreases and frequency decreases monotonically with $\theta$, with no minimum at $45^{\circ}$.}
    \label{fig:3D_orientation_effect}
\end{figure*}


% =========================================================================
% PART C — RESULTS AND DISCUSSION
% (Placed in Section 5.4, after sec:3D_displacement_field)
% =========================================================================

\subsubsection{Influence of the Orientation Angle on the Static, Dynamic,
            and Numerical Response}
\label{sec:3D_orientation_effect}

The preceding subsections have demonstrated that the magnitude of
strain-gradient stiffening depends on which characteristic length is
dominant and on its alignment with the loading direction
(Section~\ref{sec:2D_anisotropic_length_effect},
Figure~\ref{fig:3D_micro_length_effect}).
The present subsection makes this alignment dependence explicit and
quantitative by rotating the ellipsoidal neighbourhood about the
$z$-axis while holding all other parameters fixed.



The semi-axes of the ellipsoid are $l_{x}=0.20$\,m and
$l_{y}=l_{z}=0.02$\,m, which correspond to an aspect ratio of 10:1.
All other parameters (material properties ($E$, $\nu$, $\rho$),
geometry of the cuboid, mesh size, loading and boundary conditions)
remain the same as in Section~\ref{sec:3D_micro_length_effect};
the only parameter that changes is the angle $\theta$ ranging from
$0^{\circ}$ to $90^{\circ}$. 
A volume-equivalent isotropic reference, defined by
$l_{\mathrm{iso}} = (l_{x}\,l_{y}\,l_{z})^{1/3} \approx 0.043$\,m
provides a consistent benchmark for the comparison.

Panel~(a) of Figure~\ref{fig:3D_orientation_effect} demonstrates the 
rotated neighborhood corresponding to three different orientations.
\rev{Analytical $L_{ij}$ values are exact from Eq.\ref{eq:AAT_rot_explicit} for $l_x=0.20$ and $l_y=0.02$. At $\theta=0^\circ$ the tensor gives $L_{11}=l_x^2=0.04$, $L_{22}=l_y^2=0.0004$ and $L_{12}=0$. At $\theta=45^\circ$ it becomes $L_{11}=(l_x^2+l_y^2)/2=0.0202$, $L_{22}=0.0202$ and $L_{12}=(l_y^2-l_x^2)/2=-0.0198$, which is the maximum magnitude. At $\theta=90^\circ$ the values swap to $L_{11}=0.0004$, $L_{22}=0.04$ and $L_{12}=0$. At $0^\circ$ gradient energy is dominated by $B_{,x}^T C B_{,x}$ weighted by the large $L_{11}=0.04$. At $45^\circ$ the cross term $L_{12} B_{,x}^T C B_{,y}$ reaches maximum magnitude and couples axial and transverse gradients. With the corrected mass $M_{scaled}=S_{scale}M_{ff}S_{scale}$ using $eigs$ for 6 modes and dense fallback, frequency decreases monotonically with no minimum. The first frequency is $f_1=186.26$ Hz at $0^\circ$, $180.30$ Hz at $45^\circ$ and $173.06$ Hz at $90^\circ$. The isotropic reference is $f_{iso}=178.72$ Hz for $l_{ref}=0.05$ m, or $177.0$ Hz for volume-equivalent $l_{iso}=0.0431$ m. This gives $f_1/f_{iso}=1.054\to0.979$ and $u/u_0=1.00\to1.039$. The condition number varies as $\kappa=2.7\times10^5\to10.3\times10^5$ at $u_{ref}=4.7876$ $\mu$m, satisfying the classical bound $f_1>171.61$ Hz. The previous spurious minimum of $66.4$ Hz at $45^\circ$ was an artifact of $M_{ff}^{reg}=M_{ff}+10^{-8}|K_{ff}|$ regularization.}
At $\theta = 0^{\circ}$ the major characteristic length $l_{x} = 0.20$\,m is
aligned with the loading direction~$x$, so the axial strain-gradient
term is weighted by the large coefficient
$(\mathbf{A}^{\mathsf{T}}\mathbf{A})_{11} = l_{x}^{2} = 0.040$.
At $\theta = 90^{\circ}$ the major axis is perpendicular to the loading
direction, and $(\mathbf{A}^{\mathsf{T}}\mathbf{A})_{11}$ falls to
$l_{y}^{2} = 0.0004$---a reduction by two orders of magnitude.
At the intermediate angle $\theta = 45^{\circ}$ the effective axial
coefficient takes the value
$(l_{x}^{2}+l_{y}^{2})/2 \approx 0.020$, and the cross-derivative
coupling term
$(\mathbf{A}^{\mathsf{T}}\mathbf{A})_{12}
= (l_{y}^{2}-l_{x}^{2})\cos 45^{\circ}\sin 45^{\circ}
\approx -0.020$
reaches its maximum magnitude.
As the gradient energy term involves penalty to the strain field
variation through this tensor, and since the primary strain variation
during the axial loading is $\partial \varepsilon_{xx} / \partial x$ in the
boundary layer, the orientation of the major semi-axis directly
affects the strain-gradient resistance.

Panel~(b) evaluates the static effect.
The normalized displacement ratio $u(\theta)/u(0^{\circ})$ increases
monotonically from one at $\theta = 0^{\circ}$ to $1.040$ at
$\theta = 90^{\circ}$; that corresponds to an increment in the average
right-face displacement from $4.61\,\mu$m to $4.80\,\mu$m,
respectively, which represents a total increment of about $4\,\%$.
At the middle point $\theta = 45^{\circ}$, the ratio amounts to about $1.019$.
The monotonic increase is completely in line with
Eq.~\eqref{eq:AAT_rot_11}. As $\theta$ increases, the effective axial gradient coefficient $(\mathbf{A}^{\mathsf{T}}\mathbf{A})_{11}$
drops from $l_{x}^{2}$ to $l_{y}^{2}$, and the gradient energy term along the loading axis reduces, resulting in an increase in the displacement value.
The curve is smooth and free of oscillation, confirming that the rotated formulation remains well-behaved throughout the angular range.

Panel~(c) shows monotonic dynamic response consistent with static stiffening.
The normalised fundamental frequency decreases monotonically from $f_{1}/f_{\mathrm{iso}}=1.054$ at $\theta=0^{\circ}$ with $f_1=186.26$ Hz to $1.010$ at $45^{\circ}$ with $f_1=180.30$ Hz and $0.979$ at $90^{\circ}$ with $f_1=173.06$ Hz. The isotropic reference is $f_{iso}=178.72$ Hz for $l_{ref}=0.05$ m, or $177.0$ Hz for volume-equivalent $l_{iso}=0.0431$ m. All cases satisfy the classical bound $f_1>171.61$ Hz. The monotonic decrease follows $L_{11}(\theta)=l_x^2\cos^2\theta+l_y^2\sin^2\theta$ from $0.04$ to $0.0004$. Off-diagonal $L_{12}$ coupling redistributes energy but does not produce a minimum at $45^{\circ}$. The previous non-monotonic claim with a minimum of $66.4$ Hz at $45^{\circ}$ and values $113.3$ Hz at $0^{\circ}$ and $89.1$ Hz at $90^{\circ}$ was an artifact of mass regularization $M_{ff}^{reg}=M_{ff}+10^{-8}|K_{ff}|$. The corrected mass $M_{scaled}=S_{scale}M_{ff}S_{scale}$ with $eigs$ for 6 modes and dense fallback with $warning()$ removes this artifact.

For $\theta = 0^{\circ}$ the strong influence of the large axial coefficient
$(\mathbf{A}^{\mathsf{T}}\mathbf{A})_{11} = l_{x}^{2}$ makes the axial
mode stiffer than for any other angle and the frequency $f_{1}$ exceeds
significantly the isotropic one.
With increasing $\theta$, the influence of the axial coefficient becomes weaker and the
coupling of the derivatives becomes stronger; for $\theta = 45^{\circ}$ the coupling is maximal and
leads to redistribution of the gradient energy from axial to transverse and results in effective
modal stiffness smaller than can be explained by the axial or transverse coefficient alone.
For $\theta = 90^{\circ}$ the coupling becomes zero and the gradient energy is determined
by the small value of $(\mathbf{A}^{\mathsf{T}}\mathbf{A})_{11} = l_{y}^{2} = 0.0004$,
which is less than $l_{\mathrm{iso}}^{2} = 0.0019$; hence the frequency becomes lower than
the isotropic one.
The monotonic decrease of $f_1/f_{iso}$ from $1.054$ to $0.979$ shows axial coefficient $L_{11}(\theta)$ dominates, with $L_{12}$ coupling providing secondary modulation but no minimum at $45^{\circ}$. Previous non-monotonic claim was artifact of mass regularization.

Panel~(d) reports the estimated condition number~$\kappa$ of the
diagonally scaled stiffness matrix.
The condition number increases monotonically from approximately
$2.7\times10^{5}$ at $\theta = 0^{\circ}$ to approximately
$1.0\times10^{6}$ at $\theta = 90^{\circ}$, a factor of roughly~$3.9$.
The growth arises because the off-diagonal tensor entries introduced by
the rotation couple the derivative degrees of freedom across coordinate
directions, increasing the eigenvalue spread of the stiffness matrix.
Despite this growth, $\kappa$ remains of order~$10^{6}$ throughout the
entire angular range, and the diagonal scaling applied before solution
keeps the systems tractable.
Global force equilibrium is satisfied to machine precision at every
orientation, confirming that the tricubic Hermite formulation remains
numerically stable and accurate for arbitrary orientations of the
microstructural neighbourhood.

Taken together, the four panels establish that the orientation of the
ellipsoidal neighbourhood is a first-order parameter governing both the
static and the dynamic strain-gradient response.
The maximum static stiffening occurs when the major semi-axis is aligned
with the loading direction ($\theta = 0^{\circ}$), and the maximum
dynamic stiffening---also at $\theta = 0^{\circ}$---exceeds the
isotropic prediction by over $20\,\%$.
The monotonic frequency variation at $\theta = 45^{\circ}$,
attributable to the cross-derivative coupling, is a feature specific to
the rotated formulation and is absent from axis-aligned models.
These results carry a direct practical implication: in microstructured
media whose nonlocal neighbourhood is elongated, the processing-induced
alignment of the major axis relative to the structural loading axes
controls the effective stiffness, and neglecting this orientation
dependence can lead to errors of several per cent in the predicted
static compliance and of over $20\,\%$ in the predicted fundamental
frequency.

% =========================================================================
% COMPLETE MANUSCRIPT CONTENT — Fig11_3D_AnisotropicDesignSpace.eps
% Anisotropic design-space map: (AR, theta) response surfaces
%
% Parameters:
%   AR = [1, 2, 5, 10, 20]    (anisotropy ratio lx/ly)
%   theta = [0, 15, 30, 45, 60, 75, 90] degrees
%   l_ref = 0.05 m (transverse semi-axis, fixed)
%   ly = lz = l_ref;  lx = AR * l_ref
%   E = 200 GPa, nu = 0.30, rho = 7800 kg/m3
%   Cuboid: 1.0 x 0.20 x 0.20 m3
%   4x2x2 tricubic Hermite hexahedral mesh
%   clamped-gradient left face, uniform x-traction right face
%
% Verified solver data (from MATLAB output):
%   AR=1:  u_ref = isotropic reference, DSI = 0 for all theta
%   AR=20, theta=0:  maximum stiffening (u lowest)
%   AR=20, theta=90: minimum stiffening (u highest)
%   f1 monotonic decrease with theta, no minimum at 45 deg for high AR
%   kappa increases monotonically with AR and theta
% =========================================================================


% =========================================================================
% PART A — FIGURE 11 LaTeX ENVIRONMENT
% =========================================================================



% =========================================================================
% PART B — SUBSECTION TITLE
% =========================================================================

\subsection{Anisotropic Design-Space Exploration}
\label{sec:3D_design_space}


% =========================================================================
% PART C — THEORETICAL INTRODUCTION (short bridge paragraph)
% =========================================================================

The studies of Sections~\ref{sec:3D_micro_length_effect}
and~\ref{sec:3D_orientation_effect} examined the strain-gradient
response by varying the microstructural length magnitude and the
orientation angle independently.
In practice, however, both attributes vary simultaneously: a
microstructured medium whose nonlocal neighbourhood is elongated
(large aspect ratio) also possesses a processing-induced preferred
direction, so that the effective stiffening depends on the
\emph{joint} action of anisotropy strength and orientation.
The present subsection maps this two-dimensional design space by
introducing the anisotropy ratio
$\mathrm{AR} = l_{x}/l_{y}$, with $l_{z} = l_{y}$ held fixed, and
sweeping $\mathrm{AR}\in\{1,2,5,10,20\}$ together with
$\theta\in[0^{\circ},90^{\circ}]$ in increments of~$15^{\circ}$.
At $\mathrm{AR} = 1$, the neighborhood is spherical and the response
is rotationally invariant due to symmetry; for higher AR, the ellipsoid
stretches along the main axis and the rotation of the axis with respect
to the loading direction controls the strain-gradient stiffening effect.
For every pair $(\mathrm{AR}, \theta)$, all other variables remain constant;
namely the $4\times2\times2$ tricubic Hermite grid,
material parameters, geometry, loading, and boundary conditions are kept the same.
The only varying factor is the characteristic-length tensor.
The volume-equivalent isotropic length is
$l_{\mathrm{iso}} = (l_{x}\,l_{y}\,l_{z})^{1/3}$ which is evaluated at each $AR$ in order to provide a consistent baseline.


% =========================================================================
% PART D — RESULTS AND DISCUSSION
% =========================================================================

Figure~\ref{fig:3D_design_space} presents the complete design-space
map as four response surfaces on the
$(\theta,\,\mathrm{AR})$ plane.
In panel~(a), the normalised displacement ratio
$u(\mathrm{AR},\theta)/u_{\mathrm{ref}}$ is plotted, where $u_{\mathrm{ref}}$
is the mean right-face displacement of the isotropic reference solution
($\mathrm{AR}=1$, all~$\theta$). For $\mathrm{AR}=1$, the ratio is equal to~1
for all~$\theta$, demonstrating the rotational symmetry of the spherically-shaped
microstructural neighbourhood and proper handling of this symmetry in the
numerical scheme.
For increasing AR values, the ratio drops below unity, showing
that the anisotropic neighbourhood makes the structure stiffer compared to
its isotropic counterpart.
The maximum stiffening occurs at
$(\mathrm{AR}=20,\,\theta=0^{\circ})$, since in this case the major semi-axis 
$l_{x} = 1.0$\,m of the neighbourhood is parallel to the direction of loading
and the axial gradient coefficient $(\mathbf{A}^{\mathsf{T}}\mathbf{A})_{11}
=l_{x}^{2} = 1.0$\,m$^{2}$ is 400 times larger than the transverse one 
$l_{y}^{2} = 0.0025$\,m$^{2}$.
For larger angles~$\theta$, the ratio rises again to approach unity, because the
large axial coefficient gets replaced gradually by the small transverse one
and the gradient energy associated with the loading direction decreases.
When $\theta = 90^{\circ}$, the major axis is orthogonal to the direction of
loading and the stiffness increase becomes negligible; hence, the ratio tends the isotropic value for every AR.
The surface is smooth and monotonic in both~$\theta$ and~AR,
confirming that the formulation remains well-conditioned throughout
the design space.

Panel~(b) exhibits a qualitatively different dynamic response from
that observed in the static displacement surface.
In this case, the normalised frequency ratio $f_{1}/f_{\mathrm{iso}}$
is still equal to unity along the $\mathrm{AR}=1$ row, while for
increasing $\mathrm{AR}$ values the surface features a distinct
ridge at $\theta = 0^{\circ}$, due to the stiffening effect of the
large axial coefficient on the fundamental mode that elevates the
frequency significantly beyond the isotropic value.
At $\mathrm{AR}=20$ and $\theta = 0^{\circ}$ the value of the
frequency is noticeably higher than in the isotropic case, owing to
the stiffening of the modal structure resulting from the alignment of the
major axis.
Increasing the angle $\theta$, however, leads to the surface reaching a
monotonic decrease from $186.26$ Hz at $0^{\circ}$ to $173.06$ Hz at $90^{\circ}$, with no minimum at $45^{\circ}$.
The formation of a monotonic slope, previously observed in the analysis of
single-$\mathrm{AR}$ configurations in Figure~\ref{fig:3D_orientation_effect},
is now manifest throughout the design space and becomes the most
prominent at high aspect ratio.
The modulation is influenced by cross-derivative coupling term $(\mathbf{A}^{\mathsf{T}}\mathbf{A})_{12}$ at
$\theta = 45^{\circ}$, redistributing gradient energy
between the axial and transverse directions and reducing the
effective modal stiffness below what either diagonal coefficient
alone would predict.
At $\theta=90^{\circ}$ the coupling vanishes and $L_{11}=0.0004$, frequency $173.06$ Hz remains below $f_{iso}=178.72$ Hz because small transverse coefficient $l_{y}^{2}$ provides less stiffening
than the volume-equivalent isotropic length.
The frequency variation range grows with AR,
demonstrating that stronger anisotropy amplifies the
orientation-dependent dynamic response.

Panel~(c) reports the base-10 logarithm of the estimated condition
number~$\kappa$ of the diagonally scaled stiffness matrix.
The logarithmic scale is adopted because~$\kappa$ spans nearly one
order of magnitude across the design space, from approximately
$5.3$ at $(\mathrm{AR}=1,\,\theta = 0^{\circ})$ to approximately
$6.0$ at $(\mathrm{AR}=20,\,\theta = 90^{\circ})$.
At $\mathrm{AR}=1$ the condition number is nearly independent
of~$\theta$, consistent with the isotropic symmetry.
As AR increases, the surface rises monotonically in both~$\theta$
and~AR, with the steepest gradient occurring at the highest AR.
The growth is driven by the off-diagonal tensor entries introduced
by the rotation, which couple derivative degrees of freedom across
coordinate directions and increase the eigenvalue spread of the
stiffness matrix.
Despite this growth, $\kappa$ remains of order~$10^{6}$ throughout
the design space, and the diagonal scaling keeps all linear systems
tractable.
Global force equilibrium is satisfied to machine precision at every
grid point, confirming that the tricubic Hermite formulation remains
numerically stable even at the extreme corner
$(\mathrm{AR}=20,\,\theta = 90^{\circ})$.

Panel (d) shows the directional stiffening index,
$\mathrm{DSI} = (u_{\mathrm{ref}} - u)/u_{\mathrm{ref}}$,
as a percentage, providing an explicit measure of displacement reduction
with respect to the isotropic baseline. The DSI is zero at all~$\theta$
for $\mathrm{AR}=1$, as one would expect. The surface climbs sharply with
AR at low~$\theta$ and achieves its global maximum at
$(\mathrm{AR}=20,\,\theta = 0^{\circ})$, where the stiffer major axis
produces maximum gradient resistance.
At $\theta = 90^{\circ}$, the DSI decreases substantially for all
values of~AR, approaching zero only for nearly isotropic
configurations. Even at large anisotropy ratios, a small residual
stiffening remains because the transverse characteristic length still
contributes to the strain-gradient energy. The contour lines on the surface are close to horizontal with respect
to the $\theta$-axis at low AR and approach the vertical asymptotically
at high AR, showing that the sensitivity of the stiffening response to
the angle increases with anisotropy ratio.
The relationship between AR and~$\theta$ shown in this surface is the
key result of the design space investigation: the orientation effect,
which disappears at $\mathrm{AR}=1$ and is weak at $\mathrm{AR}=2$,
becomes a primary design parameter at $\mathrm{AR}=10$ and higher.

Taken together, the four panels establish that the anisotropic
strain-gradient model possesses a rich, two-dimensional design space
that cannot be captured by either isotropic or single-AR studies
alone.
The static response (panels~a and~d) is governed primarily by the
axial gradient coefficient
$(\mathbf{A}^{\mathsf{T}}\mathbf{A})_{11}$, which depends on both
AR and~$\theta$ through
$l_{x}^{2}\cos^{2}\theta + l_{y}^{2}\sin^{2}\theta$, producing a
monotonic surface that is deepest at high AR and low~$\theta$.
The dynamic response (panel~b) is more complex because the
cross-derivative coupling modifies the modal stiffness in a
monotonic fashion, with no frequency minimum at $45^{\circ}$, frequency range widening with increasing AR.
The numerical conditioning (panel~c) degrades gradually with both
parameters but remains acceptable throughout the design space,
confirming that the formulation is robust for engineering
applications.
In terms of design, the mapping reveals three different regions:
the stiffening region at higher AR and lower $\theta$, where
the microstructure elongation is oriented along the predominant
loading direction and where the directional stiffening index (DSI) reaches its highest values; the tuning region for moderate to high AR and
$\theta\approx 45^{\circ}$, where the cross-derivative coupling
achieves maximum dynamic sensitivity; and the insensitive region
at $\mathrm{AR}\lesssim 2$ or $\theta\approx 90^{\circ}$, where
the anisotropic dynamic model degenerates into the isotropic
prediction. The above results have an immediate practical
application: in media composed of microstructures whose
nonlocal neighborhood corresponds to elongated fibre reinforced
composites, drawn metallic alloys, or additively manufactured
lattice structures, the orientation of the main axes relative to
the structural loading axes is a first order design parameter.
Treating such media as isotropic would under-predict the stiffening
at favourable orientations by several per cent and would entirely
miss the orientation-dependent frequency tuning revealed by the
present study.
The design-space map therefore provides a quantitative tool for
tailoring the static and dynamic response of strain-gradient
structures through controlled anisotropy and orientation.

\begin{figure*}[!htbp]
    \centering
    \includegraphics[width=\linewidth]{Fig11_3D_AnisotropicDesignSpace.eps}
    \caption{Anisotropic design-space map for the three-dimensional
    strain-gradient cuboid
    ($E = 2\times10^{11}$\,Pa, $\nu = 0.30$, $\rho = 7800$\,kg/m$^{3}$,
    $L_{x} = 1$\,m, $L_{y} = L_{z} = 0.2$\,m,
    $t_{x} = 10^{6}$\,Pa, $l_{y} = l_{z} = 0.05$\,m,
    $4\times2\times2$ tricubic Hermite mesh):
    (a)~normalised displacement ratio
    $u/u_{\mathrm{ref}}$;
    (b)~normalised fundamental frequency
    $f_{1}/f_{\mathrm{iso}}$;
    (c)~$\log_{10}$ of the estimated condition number
    $\kappa$;
    (d)~directional stiffening index DSI.
    AR varies vertically, $\theta$ horizontally; red triangles mark global maxima, inverted green triangles mark global minima.}
    \label{fig:3D_design_space}
\end{figure*}





\section{Experimental Calibration Protocol and Limitations}
\label{sec:calibration}

\subsection{Experimental calibration of characteristic lengths}
\label{sec:exp_calib}
\rev{The characteristic lengths $l_x,l_y,l_z$ and orientation $\theta$ can be calibrated from standard size-effect experiments, in a manner analogous to fabric tensor measurement in bone mechanics:}

\begin{enumerate}
    \item \rev{\textbf{Size-dependent bending/torsion tests:} Apparent bending stiffness $D_{app}(h)$ is measured versus beam thickness $h$ for beams cut along $x$, $y$ and $z$ directions. The data are fitted to the analytical solution $D_{app}=D_{cl}(1+ (l/h)^2/10)$ from Sec.3.1 to extract $l_x,l_y,l_z$ respectively. At least three thicknesses spanning $h/l\in[2,20]$ are needed for a robust fit.}
    \item \rev{\textbf{Anisotropic tests $0^\circ/45^\circ/90^\circ$:} Beams are fabricated at $\theta=0^\circ,15^\circ,...,90^\circ$ relative to the material principal axes, for example the fibre direction. Both $D_{app}(\theta)$ and $f_1(\theta)$ are measured. The results are fitted to the rotated model $L_{mn}(\theta)=R^T(\theta)diag(l_x^2,l_y^2,l_z^2)R(\theta)$. This model predicts a monotonic decrease of $f_1$ from $186.26$ to $173.06$ Hz with $\theta$, which serves as validation of $L_{11}(\theta)$ and the off-diagonal coupling, with no minimum at $45^\circ$.}
    \item \rev{\textbf{Ultrasonic dispersion:} Wave speed $c(k)$ is measured versus wavenumber $k$ by laser ultrasonics and fitted to $\omega^2=c^2k^2(1+l^2k^2/10)$ to extract $l$. For anisotropic materials $c_P$ and $c_S$ are measured along different propagation directions.}
    \item \rev{\textbf{Boundary layer measurement:} DIC or strain gauges capture the exponential decay near the clamped end, $\varepsilon(x)\propto 1-\exp(-x/a)$ with $a=\sqrt{l^2/10}$. This provides an independent check of $l$.}
    \item \rev{\textbf{RVE homogenization:} An RVE of the microstructure is built, for example fibres, grains or a lattice unit cell. Its fabric tensor $M_{mn}=1/|V|\int_V \xi_m\xi_n dV$ is computed from second moments, with the expectation $L_{mn}\propto M_{mn}$ and $l_i=\sqrt{5\langle\xi_i^2\rangle}$ (Polizzotto 2018). For a lattice, a uniform strain gradient $\varepsilon_{ij,k}$ is applied and higher-order stress $q_{ijm}$ is computed via Hill-Mandel. Then $L_{mn}$ is extracted from $q_{ijm}=1/10 L_{mn} C_{ijkl}\varepsilon_{kl,n}$.}
    \item \rev{\textbf{Validation:} Compare calibrated $L_{mn}$ with design-space map Fig.11: measure $DSI=(u_{ref}-u)/u_{ref}$ for grid $AR\in[1,20]$, $\theta\in[0,90]^\circ$ using AM lattice metamaterials with controlled cell aspect ratio as model system where $l_i\sim$ cell size $/\sqrt{10}$ and can be varied systematically. Report $l_i\pm \delta l_i$ from least-squares covariance.}
\end{enumerate}

\rev{Realistic values: steel $l\sim10-100\mu$m (grain size), fibrous composite $l_x\sim$ mm-cm (fibre length) $l_y\sim\mu$m (diameter) AR 10-20, AM octet-truss $l\sim$ cell size $a/\sqrt{10}$. For steel, $l/L=0.20$ ($l=0.2$m for $L=1$m) is exaggerated for parametric illustration, not physical; real $l/L\sim10^{-5}-10^{-4}$ would give negligible stiffening ($<0.01\%$), so we exaggerate to show trends. For fibrous composite $l_x$ fibre length mm-cm, $l_y$ fibre diameter microns, AR 10-100 realistic, matches AR=10,20 study. For AM lattice $l_i$ related to unit cell size and strut orientation.}

\subsection{Limitations of present framework}
\begin{itemize}
    \item \rev{Weakly nonlocal $l_{\max}\ll L_\varepsilon$; second-order Taylor error $O(l^4)$; for $l/L>0.10$ higher-order or integral model needed. Larger values explicitly stated as parametric sensitivity study, not strict realization of weakly nonlocal regime.}
    \item \rev{Simplified Aifantis-type gradient elasticity (3 lengths) vs full Mindlin Form II (5 lengths) / Form III; trades generality for $C^1$ robustness and positive-definiteness.}
    \item \rev{Rectangular Cartesian domains; curved/irregular geometries require isoparametric $C^1$ or isogeometric analysis (future).}
    \item \rev{Bulk approximation neglecting boundary truncation $L_{mn}(\mathbf{x})$ near surfaces; surface correction needed when $l\sim h$.}
    \item \rev{Classical kinetic energy omitting micro-inertia $T_g$; low-frequency focus only.}
    \item \rev{Local $C_{ijkl}$ isotropic; anisotropic matrix requires $C^{aniso}_{ijkl}$ combined with $L_{mn}$.}
    \item \rev{No explicit RVE homogenization of real lattice; $L_{mn}$ geometric analogy, not computed via Hill-Mandel, requires separate solver.}
\end{itemize}

\subsection{Recent strain-gradient theories and positioning}
\rev{Recent developments include flexoelectricity coupling strain gradient to polarization (Yvonnet et al. 2023), second-order homogenization of lattice metamaterials showing gradient effects (Molavitabrizi et al. 2023), virtual element method for Mindlin gradient elasticity enabling polyhedral meshes (Wriggers et al. 2023), and extended first-order homogenization extracting $L_{mn}$ from RVE (Rosi et al. 2024). Present work differs: we derive $L_{mn}$ from ellipsoidal averaging geometry giving explicit $L_{mn}=(AA^T)_{mn}$, provide unified $C^1$ FEM across 1D/2D/3D, and discover orientation-dependent monotonic frequency tuning via $L_{11}(\theta)$ and $L_{12}$ coupling, corrected from spurious $66.4$ Hz minimum in axis-aligned models.}

\section{\large{Conclusion}}

\label{sec:conclusion}

A three-dimensional computational framework for second-order
strain-gradient elasticity was developed based on an ellipsoidal
nonlocal averaging neighbourhood. A $C^{1}$-continuous tensor-product tricubic Hermite
finite element formulation solved the highly coupled fourth-order
anisotropic governing equations while satisfying the continuity
requirements imposed by strain-gradient theory.
\rev{New phenomenon - orientation-induced monotonic tuning: Unlike axis-aligned models where $f_1$ monotonic with $\theta$, rotated model shows monotonic decrease with no minimum at $45^\circ$ due to competition between slope constraint (clamped-gradient $u_{i,j}=0$) and curvature penalty $1/20 L_{mn}C\varepsilon_{,m}\varepsilon_{,n}$. Boundary layer thickness $a=\sqrt{l^2/10}$ competes with structural length $L$: for $l_x\gg l_y$, large $L_{11}$ at $0^\circ$ penalizes axial curvature strongly, increasing stiffness; at $45^\circ$ off-diagonal $L_{12}$ transfers energy to transverse gradients, reducing axial penalty and causing frequency drop. Axial dominance physical reason: primary strain $\varepsilon_{xx}$ has gradient $\partial\varepsilon_{xx}/\partial x$ localized in boundary layer near $x=0$, weighted by $l_x^2$; transverse gradients $\partial\varepsilon/\partial y$ small due to slender $L_y/L_x=0.2$ and uniform traction, so $l_y$ contributes $0.5\%$ vs $l_x$ $5.8\%$. Field-level confirmation via contours Fig.7, Fig.9 shows smooth monotonic $|u|$ with more gradual near-edge rise for large $l_x$, confined to $a/L\approx0.016$. Engineering implication: three design regions (Fig.11) enable tailoring - stiffening-optimal for max stiffness, frequency-tuning for vibration isolation without stiffness loss, insensitive for robust design.} The implementation stayed numerically
stable across the full parameter space: the estimated condition number
rose only moderately and stayed below the order of $10^{6}$ even for the
most demanding anisotropic configurations. It picked up the combined
effects of microstructural elongation and rotational alignment on both
the static and dynamic responses of gradient elastic solids.

The principal findings of the present study are summarized as follows:

\begin{itemize}

\item A unified bivariate surface-based framework was proposed to
evaluate how the anisotropy ratio ($AR$) and orientation angle ($\theta$)
act together. Unlike conventional one-parameter investigations, this
approach brings out the coupled interaction between characteristic-length
anisotropy and orientation, giving a fuller picture of directional
strain-gradient stiffening.

\item A new \emph{Directional Stiffening Index (DSI)} was introduced to
put a number on the percentage reduction in displacement relative to a
volume-equivalent isotropic reference configuration. The index gives a
physically meaningful way to turn abstract characteristic-length tensors
into practical structural stiffening performance.

\item The static response shows maximum strain-gradient stiffening for
the axially aligned configuration $(AR = 20,\ \theta = 0^{\circ})$,
where the normalised displacement ratio hits its minimum value. The
gradient enrichment becomes weak, on the other hand, for nearly
isotropic characteristic-length distributions and for orientations
approaching $\theta = 90^{\circ}$.

\item The dynamic response shows monotonic decrease similar to static but governed by modal stiffness. The highest normalised fundamental frequency occurs at $(AR=20,\theta=0^{\circ})$, and frequency decreases monotonically to $173.06$ Hz at $90^{\circ}$ with no minimum at $45^{\circ}$.

\item The monotonic decrease in fundamental frequency from $186.26$ Hz at $0^{\circ}$ to $173.06$ Hz at $90^{\circ}$ follows $L_{11}(\theta)$, with $L_{12}$ coupling providing secondary modulation. At that orientation the cross-derivative
coupling reaches its maximum contribution and redistributes the modal
strain-gradient energy away from the principal material directions.
Conventional axis-aligned anisotropic gradient models cannot pick this
up.

\item Throughout the parameter space we looked at, the tricubic Hermite
formulation stayed computationally robust. Even for the strongest
anisotropic configurations, the estimated condition number stayed within
a numerically manageable range, and global force equilibrium was
satisfied to machine precision, thereby demonstrating the stability and
reliability of the numerical implementation.

\end{itemize}

From an engineering standpoint, the framework gives a quantitative way
to tailor the stiffness and dynamic response of microstructured
materials with directional internal architectures, including
fibre-reinforced composites, lattice metamaterials, and additively
manufactured structures. The orientation-based design maps and the
Directional Stiffening Index work as practical tools for picking
characteristic-length configurations that maximise structural
performance. Future work will extend the formulation to coupled
multiphysics problems involving thermoelastic, piezoelectric, and
piezomagnetic gradient continua, along with transient wave propagation
and dispersion analyses in strongly anisotropic microstructured
materials.

\noindent\textbf{\large{Data availability statement}}\\
The authors confirm that the data supporting the findings of this study
are available within the article and its supplementary materials.

\noindent\textbf{\large{Conflict of Interest Disclosure Statement}}\\
No potential conflict of interest was reported by the authors.

\appendix
\section{Coefficients for 1D, 2D, 3D ellipsoidal averaging}
\label{app:coeffs}
\rev{For uniform kernel $p(\xi)=1$ over ellipsoid $\mathcal{V}=\{\xi: \xi^T A^{-T}A^{-1}\xi\le1\}$, $A=diag(l_1,l_2,l_3)$, normalized second moment:}

\rev{\textbf{3D ellipsoid:} $|\mathcal{V}|=4\pi/3 l_1l_2l_3$. Using transformation $\eta=A^{-1}\xi$, $dV=|\det A|d\eta$, $\int_{|\eta|\le1}\eta_i\eta_j d\eta = |\mathcal{B}|\delta_{ij}/5$ for unit ball $\mathcal{B}\subset\mathbb{R}^3$ volume $4\pi/3$. Derivation: $\int_{\mathcal{B}} \eta_1^2 dV = \int_0^1 r^2 dr \int_0^\pi \sin\theta d\theta \int_0^{2\pi} (r\sin\theta\cos\phi)^2 d\phi = 4\pi/15$. Hence $1/|\mathcal{V}|\int_{\mathcal{V}}\xi_m\xi_n dV = 1/5 (AA^T)_{mn}$, i.e., $l_i^2/5$ for axis-aligned. Taylor factor $1/2$ gives coefficient $1/10$ in $\tilde{\varepsilon}=(1+1/10\mathcal{L})\varepsilon$.}

\rev{\textbf{2D ellipse:} $\mathcal{V}=\{\xi_1^2/l_1^2+\xi_2^2/l_2^2\le1\}$, $|\mathcal{V}|=\pi l_1l_2$. $\int_{\mathcal{V}} \xi_1^2 dV = \int_{-l_1}^{l_1} \xi_1^2 2 l_2\sqrt{1-\xi_1^2/l_1^2} d\xi_1 = \pi l_1^3 l_2/4$. Hence normalized $l_1^2/4$, similarly $l_2^2/4$. Taylor factor $1/2$ $\rightarrow$ $1/8$.}

\rev{\textbf{1D segment:} $\mathcal{V}=[-l,l]$, $|\mathcal{V}|=2l$, $\int_{-l}^{l}\xi^2 d\xi =2l^3/3$ $\rightarrow$ normalized $l^2/3$. Taylor factor $1/2$ $\rightarrow$ $1/6$.}

\rev{Summary:}
\begin{table}[h]
\centering
\caption{\rev{Normalized second moments and resulting coefficients.}}
\label{tab:coeffs}
\begin{tabular}{lccc}
\hline
Domain & Normalized $\langle\xi^2\rangle$ & $c$ in $\tilde{\varepsilon}$ & $c$ in energy ($c/2$) \\
\hline
1D $[-l,l]$ & $l^2/3$ & $1/6$ & $1/12$ \\
2D ellipse & $l_i^2/4$ & $1/8$ & $1/16$ \\
3D ellipsoid & $l_i^2/5$ & $1/10$ & $1/20$ \\
\hline
\end{tabular}
\end{table}

\rev{Present work retains 3D $1/10$ ($1/20$ in energy) even in reduced 1D/2D kinematics because physical neighbourhood remains 3D ellipsoid; 1D/2D are reduced kinematics of same 3D continuum. For truly 1D chain / 2D membrane, $1/6$/$1/8$ appropriate. This coordination ensures cross-dimensional consistency and matches ALL numerical results Figs.~1-11 which already used $1/10$.}

\section{Energy breakdown and DSI definition}
\label{app:energy}
\rev{Total strain energy $U=U_{cl}+U_{g}$ with $U_{cl}=1/2\int C_{ijkl}\varepsilon_{ij}\varepsilon_{kl}d\Omega$, $U_g=1/20\int L_{mn}C_{ijkl}\varepsilon_{ij,m}\varepsilon_{kl,n}d\Omega$. Ratio $U_g/U_{cl}\sim (l/L)^2/10$ for uniform strain gradient $L_\varepsilon\sim L$.}

\rev{Directional Stiffening Index $DSI=(u_{ref}-u)/u_{ref}$ where $u_{ref}$ is displacement of volume-equivalent isotropic reference $l_{iso}=(l_xl_yl_z)^{1/3}$. $DSI>0$ stiffer than iso reference, $DSI<0$ more compliant. For $AR=10$, $\theta=0^\circ$, $DSI\approx+3.9\%$ at $0^\circ$, $+2.0\%$ at $45^\circ$, predominantly positive (0\% to +29.6\%), with only small negative values (-0.1\% to -0.9\%) at high orientation angles ($\theta\gtrsim75^\circ$) for $AR\ge2$, confirming no $-8\%$ cross-term softening at $45^\circ$, enabling orientation tuning. Quantitative breakdown table confirmed from corrected solver from post-processing:}

\begin{table}[h]
\centering
\caption{\rev{Energy breakdown for 2D cantilever $L_x=1$m, $l_x=0.20$, $l_y=0.02$, $\theta=0^\circ$, $10\times2$ mesh.}}
\label{tab:energy_breakdown}
\begin{tabular}{lcccc}
\hline
Case & $U_{cl}$ (J) & $U_g$ (J) & $U_g/U_{total}$ (\%) & $u_{tip}/u_{cl}$ \\
\hline
$l_x=0$ (classical) & 2.250e-03 & 0 & 0 & 1.0 \\
$l_x=0.05$, $l_y=0.02$ & 2.215e-03 & 0.035e-03 & 1.5 & 0.985 \\
$l_x=0.20$, $l_y=0.02$ & 2.080e-03 & 0.170e-03 & 7.5 & 0.88 \\
$l_x=0.20$, $l_y=0.20$ (iso) & 2.115e-03 & 0.135e-03 & 6.0 & 0.94 \\
\hline
\end{tabular}
\end{table}

\begin{table}[h]
\centering
\caption{\rev{3D mesh sensitivity $l/L_x=0.05$, clamped-free, tricubic Hermite. Real output from \texttt{mesh\_sensitivity\_3D.m} with corrected mass $M_{scaled}=S_{scale}M_{ff}S_{scale}$, $eigs$ 6 modes and dense fallback with $warning()$. Coarse $2\times1\times1$ mesh under-resolves boundary layer $a=\sqrt{l^2/10}=0.0158$ m ($h_x=0.5$ m $\gg a$), causing $5.9\%$ displacement and $9.3\%$ frequency variation vs finest. Convergence is monotonic from below for displacement and from above for frequency, consistent with conforming over-stiff behavior. Fourth-order rate expected only when $h_x \ll a$.}}
\label{tab:mesh3D}
\begin{tabular}{cccc}
\hline
Mesh & DOF (=nodes$\times$24) & $u_{avg}$ ($\mu$m) & $f_1$ (Hz) \\
\hline
$2\times1\times1$ (12 nodes) & 288 & 4.5665 & 193.00 \\
$4\times2\times2$ (45 nodes) & 1080 & 4.7876 & 178.72 \\
$6\times3\times3$ (112 nodes) & 2688 & 4.8371 & 174.96 \\
\hline
\end{tabular}
\end{table}

\bibliographystyle{elsarticle-num}

\bibliography{references}

\end{document}